\pdfoutput=1
\documentclass[11pt]{amsart}
\usepackage{amsmath,amssymb,amsthm}
\usepackage[margin=1.2in]{geometry}
\usepackage[colorlinks=true,linkcolor=blue,citecolor=blue,urlcolor=blue]{hyperref}
\hypersetup{pdfauthor={Bernd Johannes Wuebben},
 pdftitle={Convergence of the minimizing flow on acyclic quivers},
 pdfsubject={MSC 2020: 16G20, 53D20; 14F08, 14L24, 18G80},
 pdfkeywords={quiver representation, complex of representations, stability condition, harmonic metric, Harder-Narasimhan filtration, Kempf-Ness function, moment map, derived category, projective line}}
\setlength{\emergencystretch}{2em}
\newtheorem{theorem}{Theorem}[section]
\newtheorem{proposition}[theorem]{Proposition}
\newtheorem{lemma}[theorem]{Lemma}
\newtheorem{corollary}[theorem]{Corollary}
\theoremstyle{definition}
\newtheorem{definition}[theorem]{Definition}
\theoremstyle{remark}
\newtheorem{remark}[theorem]{Remark}
\newcommand{\CC}{\mathbb C}
\newcommand{\RR}{\mathbb R}
\newcommand{\ii}{\mathrm i}
\newcommand{\cK}{\mathcal K}
\newcommand{\cN}{\mathcal N}
\newcommand{\cS}{\mathcal S}
\newcommand{\Eone}{\mathrm E_1}
\newcommand{\mass}{\mathfrak m}
\DeclareMathOperator{\tr}{tr}
\DeclareMathOperator{\Rea}{Re}
\DeclareMathOperator{\Ima}{Im}
\DeclareMathOperator{\diag}{diag}
\DeclareMathOperator{\Hom}{Hom}
\newcommand{\norm}[1]{\left\lVert#1\right\rVert}
\newcommand{\inner}[2]{\left\langle#1,#2\right\rangle}


\title[The minimizing flow on acyclic quivers]{Convergence of the
minimizing flow\\ on acyclic quivers}
\author{Bernd Johannes Wuebben}
\email{wuebben@gmail.com}
\subjclass[2020]{Primary 16G20, 53D20; Secondary 14F08, 14L24, 18G80}
\keywords{Quiver representation, complex of representations, stability
condition, harmonic metric, Harder--Narasimhan filtration, Kempf--Ness
function, moment map, derived category, projective line}
\date{November 18, 2026\\This version: November 19, 2026}

\begin{document}
\begin{abstract}
The minimizing flow of Haiden, Katzarkov, Kontsevich and Pandit recovers
the Harder--Narasimhan filtration of every nonzero representation of a
finite acyclic quiver from the exponential decay rates of its Hermitian
metric $h(t)$. This holds for all vertex charges in the upper half-plane and every
initial metric. If $v$ lies in the $k$th filtration step outside the
preceding one and $Z_k$ is the central charge of the $k$th quotient, then
$t^{-1}\log h(t)(v,v)$ tends to $-\arg Z_k$. The mass of the evolving metric tends to
$\sum_k|Z_k|$. On each quotient the metric rescaled by $e^{t\arg Z_k}$
converges exponentially to a harmonic metric if the quotient is polystable.
In general the rescaled metric stays at bounded distance from a
Kempf--Ness gradient trajectory with vertex weights proportional to the
imaginary parts of the charges. The iterated-logarithm asymptotics of
Haiden, Katzarkov, Kontsevich and Pandit therefore apply up to bounded
factors. For two-term complexes with nonzero cohomology concentrated in
degree zero and differential supported at one vertex, the same results
hold on cohomology when the phase is taken from the polar decomposition.
For generalized Kronecker quivers the results extend to bounded complexes
whose differentials vanish at one vertex, with the phases shifted by
$-\pi n$ in degree $n$.
Under Beilinson's equivalence this includes every object of
$D^b(\mathbb P^1)$ in a minimal projective or injective Kronecker model.
\end{abstract}
\maketitle

\section{Introduction}\label{sec:intro}

Haiden, Katzarkov, Kontsevich and Pandit \cite{hkkp26} propose a notion of
K\"ahler geometry on a triangulated category with a stability condition in
the sense of Bridgeland \cite{bri07}, in which an object carries a space of
metrics, an endomorphism $\Omega$ depending on the metric whose trace is the
central charge, a phase $\arg\Omega$ defined by a positivity condition, a
mass $\tr|\Omega|$, and a flow $h^{-1}\dot h=-\arg\Omega$ of the metrics,
called the minimizing flow in \cite[\S4.1]{hkkp26}. A metric is
\emph{harmonic} if $\arg\Omega$ is a scalar, and harmonic metrics are
expected to exist on polystable objects and to be the limits of the flow
\cite[\S\S1,~4.1]{hkkp26}, a categorical form of the theorem of Donaldson,
Uhlenbeck and Yau \cite{don85,uy86}. This paper proves that on a quiver
without oriented cycles the flow converges, up to a scalar function of the
time, on every polystable representation, and that for complexes of
representations with one differential the metrics induced on the cohomology
converge in the same sense when the cohomology is polystable. When it is not
polystable, the long-time behaviour is described through the
Harder--Narasimhan filtration of the cohomology and, on its factors, through
gradient flows of Kempf--Ness functions.

\subsection{Representations}
Let $Q$ be a finite quiver without oriented cycles, with central charges
$z_i$ in the upper half-plane at its vertices, and let $\mathsf R$ be a
representation of $Q$ with Hermitian metrics $h_i$ on its spaces. Write
$\beta^\dagger$ for the adjoint of a linear map for the metrics. The
\emph{moment map} at the vertex $i$ is the Hermitian endomorphism $\mu_i$
equal to the sum of $\beta^\dagger\beta$ over the arrows $\beta$ out of $i$
minus the sum of $\beta\beta^\dagger$ over the arrows into $i$. For a
representation the endomorphism $\Omega$ of \cite[\S5.1]{hkkp26} is
$\bigoplus_i(z_i+\mu_i)$; it is normal, and the flow is
\[
 h_i^{-1}\dot h_i=-\arg(z_i+\mu_i),
\]
with the argument taken by the functional calculus, in $(0,\pi)$. Let
$Z(\mathsf R)=\sum_iz_i\dim\mathsf R_i=|Z|e^{\ii\varphi}$ be the central
charge of $\mathsf R$, and
\[
 \theta_i=\frac{\Ima(e^{-\ii\varphi}z_i)}{\sin\varphi}.
\]
A representation with central charge of argument $\varphi$ is
\emph{semistable} if no subrepresentation has a central charge of argument
larger than $\varphi$; this is King's stability \cite{kin94} for the real
parameter $-\theta$, and a metric is harmonic exactly when
$\mu_i=\theta_i$ for all $i$, that is, when it is a zero of the moment map
shifted by the stability parameter (Section~\ref{sec:convergence}).

\begin{theorem}[Representations]\label{thm:introR}
For all charges in the upper half-plane, every nonzero representation
$\mathsf R$ of $Q$ and every initial metric, along the flow:
\begin{enumerate}
\item[(i)] the Kempf--Ness function
 $\sum_a\norm{\rho_a}^2+\sum_i\theta_i\log\det h_i$ of $\mathsf R$ does not
 increase;
\item[(ii)] if $\mathsf R$ is semistable, then $\mu_i(t)\to\theta_i$ for all
 $i$, and the mass tends to $|Z|$;
\item[(iii)] if $\mathsf R$ is polystable, the metrics $e^{\varphi t}h(t)$
 converge exponentially to a harmonic metric.
\end{enumerate}
\end{theorem}

Here $\rho_a$ is the map of the arrow $a$ and $\norm\cdot$ the
Hilbert--Schmidt norm for the metrics. The factor $e^{\varphi t}$ is needed
because a harmonic metric is not stationary: it moves by the scalar
$e^{-\varphi t}$. After multiplication of the metrics by $e^{\varphi t}$ the
velocity of the flow at the vertex $i$ is $\varphi-\arg(z_i+\mu_i)$, an
increasing function of $\mu_i-\theta_i$ that vanishes at zero. The downward
gradient flow of the Kempf--Ness function, which \cite[(5.5)]{hkkp26}
describes as a simplification of this flow, has the velocity
$\mu_i-\theta_i$ up to a positive factor; its exponential convergence on
polystable representations is classical (see \cite[\S5.1]{hkkp23}, and
\cite{grs21} for the general theory), and the gradient flow of the norm
square of the moment map on the space of representations is studied in
\cite{hw11}. Theorem~\ref{thm:introR} concerns the flow $-\arg\Omega$ itself.
The precise statement, with the rate, is Theorem~\ref{thm:representations}.

\subsection{Complexes with one differential}
For complexes of representations $\Omega$ is defined through a
superconnection \cite[\S5.4]{hkkp26}; it is then not normal in general, and
we take $\arg\Omega$ to be the argument of the unitary factor in the polar
decomposition, which is our reading of the positivity condition of
\cite[\S4.2]{hkkp26}; it agrees with the functional calculus of
\cite[\S5.4]{hkkp26} when $\Omega$ is normal (Section~\ref{sec:setting}). All
statements below refer to this choice, and we call the flow the \emph{phase
flow}.

Fix a vertex $i_0$. A complex of representations $\mathsf E$ of $Q$ has
\emph{a surjective differential at $i_0$} if at each vertex $i\neq i_0$ it is
a vector space in degree zero and at $i_0$ it is $V\xrightarrow{D}U$ in
degrees $0,1$ with $D$ surjective (Definition~\ref{def:type}). Its cohomology
is a representation $\mathcal H$ of $Q$ in degree zero, with $\ker D$ at
$i_0$; we assume $\mathcal H\neq0$, and put
$Z=Z(\mathcal H)=|Z|e^{\ii\varphi}$ and $\theta$ as above. For Hermitian
metrics $h$ on the spaces of $\mathsf E$, the endomorphism $\Omega$ preserves
each of these spaces. We write $\Omega_V$, $\Omega_U$ and $\Omega_i$
($i\neq i_0$) for its blocks,
$\mass(h)=\tr|\Omega_V|+\tr|\Omega_U|+\sum_{i\neq i_0}\tr|\Omega_i|$ for the
mass, where $|T|=(TT^\dagger)^{1/2}$, and $\mu_i$ for the moment map of
$\mathcal H$ with the induced metrics, so that the endomorphism $\Omega$ of
the representation $\mathcal H$ is $\bigoplus_i(z_i+\mu_i)$. Let $\lambda$ be
the least eigenvalue of $DD^\dagger$ and $\norm\cdot_1$ the trace norm. The
quotient complex $\mathsf E/\mathcal H$ is supported at $i_0$ and is exact.
We call its metric realization $(\ker D)^\perp\to U$ the \emph{exact part}.

\begin{theorem}[Convergence on the cohomology]\label{thm:introB}
Let $\mathsf E$ be a complex with a surjective differential at $i_0$. For all
charges in the upper half-plane and every initial metric, along the phase
flow:
\begin{enumerate}
\item[(i)] if $\mathcal H$ is semistable, then $\mu_i(t)\to\theta_i$ for all
 $i$ and $\mass(h(t))\to|Z|$;
\item[(ii)] if $\mathcal H$ is polystable, the metrics induced on the spaces
 of $\mathcal H$, multiplied by $e^{\varphi t}$, converge exponentially to a
 harmonic metric of $\mathcal H$.
\end{enumerate}
\end{theorem}

\begin{theorem}[Decay of the exact part]\label{thm:introA}
Let $\mathsf E$ be a complex with a surjective differential at $i_0$ and
$U\neq0$. For all charges in the upper half-plane and every initial metric,
along the phase flow, $\lambda(t)\to\infty$ and, for $t\geq t_1$,
\[
 \norm{\Omega_V-(z_{i_0}+\mu_{i_0})\oplus0}_1+\norm{\Omega_U}_1
 +\sum_{i\neq i_0}\norm{\Omega_i-(z_i+\mu_i)}_1
 \leq C\,e^{-\varkappa(t-t_1)},
\]
where $\oplus\,0$ refers to $V=\ker D\oplus(\ker D)^\perp$, the constants
$C,\varkappa>0$ depend only on the quiver, the vertex $i_0$, the charges and
the dimensions, and $t_1$ depends in addition on the initial metric.
\end{theorem}

Theorem~\ref{thm:introA} makes no assumption on $\mathcal H$. For $U\neq0$
the complex itself carries no harmonic metric, since the phase has its
spectrum in $(0,\pi)$ on $V$ and in $(-\pi,0)$ on $U$; what is proved is that
the exact part decays and that the metrics induced on the cohomology
converge. By duality \cite[Theorem~5.1]{wue26c} the two theorems hold for
complexes \emph{with an injective differential at $i_0$}, which at $i_0$ are
$U'\to V'$ in degrees $-1,0$ with an injective differential
(Proposition~\ref{prop:injective}). The precise statements, with the
constants, are Theorems~\ref{thm:decay} and~\ref{thm:convergence},
Proposition~\ref{prop:sink} and Remark~\ref{rem:vanishing}.

\subsection{The method}
For a complex with $U\neq0$ the Kempf--Ness function of the cohomology
decreases only up to an error controlled by the exact part
(Theorem~\ref{thm:convergence}(i)), so the proof cannot start from a
Lyapunov function; it rests on an a priori bound for the arrows. Once the arrows are
bounded, the exact part decays, the error becomes exponentially small, and
the convexity of the Kempf--Ness function gives the convergence
(Proposition~\ref{prop:perturbed}). For an arrow $a$ with map $\rho_a$ let
$e_a=\norm{\rho_a}^2$ if $a$ does not start at $i_0$ and
$e_a=\tr(\rho_aR(D^\dagger D)\rho_a^\dagger)$, with $R(b)=(1-e^{-b})/b$, if it
does, and put $\cN=\sum_ae_a$. For a representation $\cN$ is the sum of the
squared norms of the arrows.

\begin{theorem}[A priori bound]\label{thm:introC}
Let $\mathsf E$ be a complex with a surjective differential at $i_0$. There
are constants $\gamma>0$ and $c\geq0$, depending only on the quiver, the
vertex $i_0$, the charges and the dimensions, such that for all metrics
\[
 \frac{d\cN}{dt}\leq-\gamma\,\cN+c
\]
along the phase flow. For a representation one can take $\gamma=\pi/\ell$,
with $\ell$ the number of arrows of a longest path in $Q$, taking $\ell=1$
when $Q$ has no arrows.
\end{theorem}

The derivative $d\cN/dt$ is a sum of three traces, one for $V$, one for $U$
and one for the remaining vertices (Lemma~\ref{lem:energyderivative}). The
trace for $V$ is bounded by a trace inequality for the polar phase, and the
trace for $U$ by a bound for the negative part of a Hermitian matrix that
occurs in $\Omega_U$, both from \cite{wue26b}
(Proposition~\ref{prop:used}(ii),(i)); the third by an elementary estimate
(Lemma~\ref{lem:three}(c)), in which the blocks are normal and are
indefinite perturbations of the charges. The estimates at the different
vertices are combined with coefficients that increase along the arrows, and
$\gamma$ is the least increment; this is where the absence of oriented
cycles is used, and for a quiver with a loop no such bound holds
(Remark~\ref{rem:energy}). See Theorem~\ref{thm:energy} and
Proposition~\ref{prop:reparrows}.

\subsection{Cohomology that is not polystable}
Every nonzero representation $\mathcal H$ has a unique filtration
$0=\mathcal H^0\subset\mathcal H^1\subset\dots\subset\mathcal H^{\mathsf n}=\mathcal H$
by subrepresentations whose factors
$\operatorname{gr}^k=\mathcal H^k/\mathcal H^{k-1}$ are semistable with
$\arg Z(\operatorname{gr}^1)>\dots>\arg Z(\operatorname{gr}^{\mathsf n})$,
the Harder--Narasimhan filtration; we write
$\phi_k=\arg Z(\operatorname{gr}^k)$. The sum
$\sum_k|Z(\operatorname{gr}^k)|$ is the mass of $\mathcal H$ for the stability
condition.

\begin{theorem}[Harder--Narasimhan asymptotics]\label{thm:introH}
Let $\mathcal H$ be a nonzero representation of $Q$ with the phase flow from
an arbitrary initial metric, or the cohomology of a complex with a surjective
or an injective differential at one vertex, with the metrics $h_i(t)$ induced
by the phase flow of the complex. For all charges in the upper half-plane:
\begin{enumerate}
\item[(i)] for every nonzero vector $v$ of $\mathcal H_i$,
 $\lim_{t\to\infty}\frac1t\log h_i(t)(v,v)=-\phi_k$, where $k$ is the least
 index with $v\in\mathcal H^k_i$; so the Harder--Narasimhan filtration is the
 filtration of the spaces by the exponential rate of decay of the metric;
\item[(ii)] the eigenvalues of the argument of $z_i+\mu_i$, with $\mu$ the
 moment map of $\mathcal H$ for the metrics $h(t)$, converge to the $\phi_k$
 with the multiplicities $\dim\operatorname{gr}^k_i$, and
 $\sum_i\tr|z_i+\mu_i|$ converges to $\sum_k|Z(\operatorname{gr}^k)|$, which
 is its infimum over all metrics on $\mathcal H$; along the phase flow of a
 representation this sum, the mass, is nonincreasing;
\item[(iii)] the second fundamental forms of the subrepresentations
 $\mathcal H^k$ tend to zero exponentially, and the metric induced on
 $\operatorname{gr}^k$, multiplied by $e^{\phi_kt}$, has the properties
 stated in Theorem~\ref{thm:introB} and in Theorem~\ref{thm:introS} below
 for $\operatorname{gr}^k$ in the place of $\mathcal H$: its moment map
 tends to the stability parameter of $\operatorname{gr}^k$, and it converges
 exponentially to a harmonic metric if $\operatorname{gr}^k$ is polystable.
\end{enumerate}
\end{theorem}

\begin{theorem}[Semistable cohomology]\label{thm:introS}
Let $\mathcal H$ be as in Theorem~\ref{thm:introH} and semistable, with
$Z(\mathcal H)=|Z|e^{\ii\varphi}$. Then the metrics $e^{\varphi t}h(t)$ remain
at bounded distance, in the space of metrics, from a trajectory of the
gradient flow of the Kempf--Ness function of $\mathcal H$ for the
Riemannian metric $\sum_i\mathsf m_i\tr\xi_i^2$ with the weights
$\mathsf m_i=\Ima z_i/\sin^2\varphi$. Hence they are an asymptotic solution of
this gradient flow in the sense of \cite{hkkp23}, and by
\cite[Corollary~5.10 and Theorem~5.11]{hkkp23} their growth and decay are
given, up to bounded factors, by products of powers of iterated logarithms
of $t$ on the steps of the iterated weight filtration of $\mathcal H$, for
the weight function $\sum_i\mathsf m_i\dim\mathsf R_i$, which equals
$|Z(\mathsf R)|/\sin\varphi$ on the subquotients $\mathsf R$ of $\mathcal H$
of phase $\varphi$.
\end{theorem}

For a nonzero arrow $\CC\to\CC$ and two charges of equal argument, the
metrics grow and decay like powers of $t$ with exponents $|z_1|/|z_0+z_1|$
and $-|z_0|/|z_0+z_1|$. These exponents determine the ratio of the two
weights (Proposition~\ref{prop:weights}). The proof of
Theorem~\ref{thm:introS} compares the velocity of the flow with the gradient
of the Kempf--Ness function: the difference is quadratic in the shifted
moment map, hence integrable in time, and gradient flows of convex
functions on spaces of nonpositive curvature do not increase distances. The
proof of Theorem~\ref{thm:introH} uses the mass as a Lyapunov function; its
dissipation is bounded below by a multiple of the sum of the squared norms
of the commutators of the arrows with the phases. See Theorems~\ref{thm:comparison} and~\ref{thm:hn},
Corollary~\ref{cor:iterated} and
Propositions~\ref{prop:dissipation} and~\ref{prop:hnmass}.

\subsection{Kronecker quivers and the projective line}
Let $K_m$ be the generalized Kronecker quiver, with two vertices and $m$
arrows from the first to the second.

\begin{theorem}[Differentials at one vertex]\label{thm:introD}
Let $\mathsf E$ be a bounded complex of representations of $K_m$ whose
differentials vanish at one of the two vertices, with any Hermitian metric.
\begin{enumerate}
\item[(i)] $\mathsf E$ is an orthogonal direct sum of subcomplexes
 $\mathsf E_n$, one for each degree $n$, where $\mathsf E_n$ has cohomology
 $H^n(\mathsf E)$ in degree $n$ and, if $H^n(\mathsf E)\neq0$, is a shift of a
 complex with a surjective or with an injective differential at one vertex.
 The phase flow preserves this decomposition and is the product of the flows
 of the $\mathsf E_n$.
\item[(ii)] For all charges in the upper half-plane and every $n$, the
 blocks of the endomorphism $\Omega_{\mathsf E_n}$ of $\mathsf E_n$ approach
 those of $(-1)^n\Omega_{H^n(\mathsf E)}$, the endomorphism of the
 representation $H^n(\mathsf E)$ with its induced metrics, extended by zero;
 the trace-norm distance decays exponentially from a time that depends on
 the initial metric. On nonzero $H^n(\mathsf E)$ the induced metrics
 multiplied by $e^{-\pi nt}$ satisfy the conclusions of
 Theorem~\ref{thm:introH} and, in the semistable case, of
 Theorem~\ref{thm:introS}. On a Harder--Narasimhan quotient of phase
 $\phi$, the induced metric multiplied by $e^{(\phi-\pi n)t}$ converges
 exponentially to a harmonic metric if the quotient is polystable. The mass
 $\mass(h(t))$ converges to the sum over $n$ of the masses of the
 representations $H^n(\mathsf E)$ for the stability condition.
\item[(iii)] Minimal bounded complexes of projective representations of
 $K_m$, and minimal bounded complexes of injective representations, have
 differentials vanishing at one vertex. Hence (i) and (ii) hold for every
 object of the bounded derived category of $K_2$, that is, of the projective
 line, represented by such a complex with arbitrary metrics; the limit of
 the mass is the mass of the object for the stability condition with heart
 the category of representations.
\end{enumerate}
\end{theorem}

See Theorems~\ref{thm:decomposition}, \ref{thm:longtime} and~\ref{thm:P1}
and Proposition~\ref{prop:minimal}.

\subsection{Relation to earlier work and scope}
For $Q=K_m$ and a surjective differential at the first vertex, with $U\neq0$
and cohomology that is nonzero at both vertices,
Theorems~\ref{thm:introB}, \ref{thm:introA} and~\ref{thm:introC} are
\cite[Theorems~5.4, 5.7, 6.1 and~6.6]{wue26b}, with the same proofs, up to
two differences: here equal arguments of the two charges are included in
Theorem~\ref{thm:introB}, and the constant $c$ of
Theorem~\ref{thm:introC} is the one of \cite{wue26b} only after the
refinement of Remark~\ref{rem:energy}(2). Sections~\ref{sec:decay}
and~\ref{sec:convergence} follow \cite{wue26b}. What is added here is the
a priori bound on a general quiver (Section~\ref{sec:energy}), a reduction of
the Kempf--Ness theorem for real stability parameters to the rational case
of King (Lemma~\ref{lem:king}), the description of the flow for cohomology
that is not polystable (Section~\ref{sec:asymptotics}), and
Section~\ref{sec:kronecker}.
Theorem~\ref{thm:introR} is contained in \cite{wue26b} only for
representations of $K_m$ with both spaces nonzero and the charge of the
first vertex of larger argument than that of the second, where it follows
from the theorems on complexes by adding an exact summand.

Complexes with nonzero differentials at two or more vertices are open, in
particular complexes of representations of $K_m$ with differentials at both
vertices. Quivers with oriented cycles are excluded. For a cohomology
representation that is semistable and not polystable the induced metrics
have no limit after any normalization by a scalar function of the time
(Remark~\ref{rem:scope}); Theorem~\ref{thm:introS} determines them up to a
bounded factor, and whether this factor converges is open in general. It
converges in the example of Proposition~\ref{prop:weights} and for a
representation of $K_2$ for which the metrics are shown in
\cite[Proposition~6.9]{wue26b} to grow and decay like $t^{\pm1/2}$. The time $t_1$ of Theorem~\ref{thm:introA} is explicit
but very large (Remark~\ref{rem:decay}).

\subsection{Organization}
Section~\ref{sec:setting} fixes the setting and states the results used from
\cite{wue26a,wue26b,wue26c}. Section~\ref{sec:energy} proves
Theorem~\ref{thm:introC}, Section~\ref{sec:decay} Theorem~\ref{thm:introA},
and Section~\ref{sec:convergence} Theorems~\ref{thm:introR}
and~\ref{thm:introB}, including the Kempf--Ness theorem for real stability
parameters in the form needed. Section~\ref{sec:asymptotics} proves
Theorems~\ref{thm:introH} and~\ref{thm:introS}, and
Section~\ref{sec:kronecker} proves Theorem~\ref{thm:introD}.

\section{Setting and results used}\label{sec:setting}

\subsection{Complexes with one differential}
Let $Q$ be a finite quiver without oriented cycles, with set of vertices
$Q_0$ and set of arrows $Q_1$; for an arrow $a$ we write $\mathsf s(a)$ and
$\mathsf t(a)$ for its source and target, and $a:i\to j$ if
$\mathsf s(a)=i$, $\mathsf t(a)=j$. Fix \emph{central charges}
$z_i=x_i+\ii y_i$, $i\in Q_0$, in the open upper half-plane, and put
$\alpha_i=\arg z_i\in(0,\pi)$. Phases are arguments, measured in radians; the
phase of \cite{bri07} is the argument divided by $\pi$, and a shift by one
degree changes arguments by $\pi$. A \emph{complex of representations} of $Q$
is a bounded complex in the abelian category of finite-dimensional
representations: vector spaces $E_i^n$, differentials
$d_i:E_i^n\to E_i^{n+1}$ and, for each arrow $a:i\to j$, a chain map
$\rho_a:E_i^\bullet\to E_j^\bullet$. For a representation $\mathsf R$ with
dimension vector $(n_i)$ we write $Z(\mathsf R)=\sum_in_iz_i$ for its central
charge.

\begin{definition}\label{def:type}
Let $i_0\in Q_0$. A complex of representations $\mathsf E$ of $Q$ is a
\emph{complex with a surjective differential at $i_0$} if at every vertex
$i\neq i_0$ it is a vector space $W_i$ in degree zero, and at $i_0$ it is
\[
 V\xrightarrow{\ D\ }U\qquad\text{in degrees }0,1,
\]
with $D$ surjective. It is a \emph{complex with an injective differential at
$i_0$} if at every vertex $i\neq i_0$ it is a vector space in degree zero,
and at $i_0$ it is $U'\xrightarrow{D'}V'$ in degrees $-1,0$ with $D'$
injective. In both cases the cohomology is a representation in degree zero,
and we assume that it is not zero.
\end{definition}

So $\mathsf E$ has cohomology in degree zero and one differential, located
at $i_0$, from degree $0$ to degree $1$ or from degree $-1$ to degree $0$.
For a complex with a surjective differential the
arrows act by linear maps $\rho_a$ between the spaces of degree zero, with
$D\rho_a=0$ for the arrows into $i_0$ and no condition on the arrows out of
$i_0$. Its cohomology is the representation $\mathcal H$ with
$\mathcal H_{i_0}=\ker D$, $\mathcal H_i=W_i$ for $i\neq i_0$ and the
restricted arrows. Put $r=\dim U\geq0$, $n_i=\dim\mathcal H_i$,
\[
 Z=Z(\mathcal H)=|Z|e^{\ii\varphi},\qquad\varphi\in(0,\pi),\qquad\sigma=\sin\varphi ,
\]
and write $z_\ast=x_\ast+\ii y_\ast$ and $\alpha_\ast$ for the charge at
$i_0$. We call $\varphi$ the \emph{total phase} of $\mathsf E$; $\varphi/\pi$
is the phase of $\mathcal H$ for the stability condition with central charge
$Z(\,\cdot\,)$ and heart the category of representations
\cite[Propositions~2.4 and~5.3]{bri07}. For
$r=0$ the complex is a representation. For the generalized Kronecker quiver
$K_m$, with vertices $0,1$ and $m\geq1$ arrows from $0$ to $1$, the complexes
with a surjective differential at the vertex $0$ with $r\geq1$ and
$n_0,n_1\geq1$ are those of \cite{wue26a,wue26b}.
Sections~\ref{sec:energy}--\ref{sec:convergence} treat complexes with a
surjective differential; those with an injective differential are their
duals (Proposition~\ref{prop:injective}).

A \emph{metric} on $\mathsf E$ consists of Hermitian metrics on $V$, $U$ and
the $W_i$; adjoints for the metrics are written ${}^\dagger$. The subcomplex
$(\ker D)^\perp\to U$ at $i_0$, which depends on the metric and is exact, is
called the \emph{exact part} of the metrized complex, and $r$ its
\emph{rank}. Put
\[
 B=D^\dagger D,\qquad L=DD^\dagger,\qquad\lambda=\lambda_{\min}(L),\qquad
 R(b)=\frac{1-e^{-b}}b,
\]
\[
 C_{\mathrm{out}}=\sum_{a:i_0\to\cdot}\rho_a^\dagger\rho_a,\qquad
 G_{\mathrm{in}}=\sum_{a:\cdot\to i_0}\rho_a\rho_a^\dagger,
\]
and, for Hermitian $\Delta\geq0$ and matrices $G$, $M$,
\[
 K_\Delta(G)=\int_0^1e^{-(1-x)\Delta}Ge^{-x\Delta}dx,\qquad
 J_L(M)=\int_{\substack{x,y\geq0\\x+y\leq1}}e^{-xL}Me^{-yL}\,dx\,dy,
\]
and $J=J_L(DC_{\mathrm{out}}D^\dagger)$. We write $M_\pm\geq0$ for the positive
and negative parts of a Hermitian matrix $M$, $\norm\cdot_1$ for the trace
norm, $\norm\cdot$ for the Hilbert--Schmidt norm and
$\norm\cdot_{\mathrm{op}}$ for the operator norm of matrices.

\subsection{The endomorphism \texorpdfstring{$\Omega$}{Omega}, the phase and the flow}
For a bounded complex of representations of $Q$ with Hermitian metrics on all
its spaces $E_i^n$, the endomorphism $\Omega$ of \cite[\S5.4]{hkkp26} is
defined through a superconnection in a graded algebra attached to $Q$:
$\Omega$ is the grading times the component along the basis elements $f_i$
of $\exp(-\sqrt{-1}F)$, where
$F=(A-A^*)^2+\sqrt{-1}\sum_if_i\otimes z_i$ and
$A=\sum_ie_i\otimes d_i+\sum_aa\otimes\rho_a$ in the algebra of
\cite[Definition~5.13]{hkkp26} tensored with the endomorphisms. Expanding the
exponential in the terms that contain the arrows, which stops at second
order, gives the following formula \cite[Proposition~2.1]{wue26c}, which may
be taken as the definition of $\Omega$ here. Let $\varepsilon$ be the grading, equal
to $(-1)^n$ in degree $n$, and for $i\in Q_0$ and an arrow $a:i\to j$ put,
on $E_i=\bigoplus_nE_i^n$,
\[
 \Delta_i=d_id_i^\dagger+d_i^\dagger d_i,\qquad
 \eta_a=d_j^\dagger\rho_a-\rho_ad_i^\dagger,\qquad
 \mu_i^{\mathsf E}=\sum_{a:i\to\cdot}\rho_a^\dagger\rho_a-\sum_{a:\cdot\to i}\rho_a\rho_a^\dagger .
\]
Then $\Omega$ preserves each $E_i^n$, and on $E_i$
\begin{equation}\label{eq:general}
 \begin{aligned}
 \varepsilon\,\Omega=K_{\Delta_i}\bigl(z_i+\mu_i^{\mathsf E}\bigr)
 &+\sum_{a:k\to i}\int e^{-x\Delta_i}\eta_a\,e^{-(1-x-y)\Delta_k}\eta_a^\dagger\,e^{-y\Delta_i}\\
 &-\sum_{a:i\to j}\int e^{-x\Delta_i}\eta_a^\dagger\,e^{-(1-x-y)\Delta_j}\eta_a\,e^{-y\Delta_i},
 \end{aligned}
\end{equation}
the integrals over $\{x,y\geq0,\ x+y\leq1\}$ with the measure $dx\,dy$. For a
representation, that is, a complex in degree zero, this is
$\Omega=\bigoplus_i(z_i+\mu_i^{\mathsf E})$: the charge at each vertex shifted
by the moment map. All terms other than
$z_ie^{-\Delta_i}$ are Hermitian, so $(-1)^n\Omega$ has the positive definite
imaginary part $y_ie^{-\Delta_i}$ on $E_i^n$. A matrix $T$ with positive
definite imaginary part has a unique Hermitian matrix $\Psi$ with spectrum in
$(0,\pi)$ such that $Te^{-\ii\Psi}$ is positive definite: in the polar
decomposition $T=|T|\,\mathsf W$, with $|T|=(TT^\dagger)^{1/2}$ and
$\mathsf W$ unitary, an eigenvalue $e^{\ii\psi}$ of $\mathsf W$ with unit
eigenvector $v$ gives $\inner v{Tv}=e^{\ii\psi}\inner v{|T|v}$, hence
$\sin\psi>0$ (see \cite[Lemma~2.2]{wue26a}). We call $\Psi$ the \emph{polar
phase} of $T$ and $|T|=Te^{-\ii\Psi}$ its modulus. The \emph{phase} $\Phi$ of
the metrized complex is, on $E_i^n$, the polar phase of $(-1)^n\Omega$ minus
$\pi n$, with spectrum in $(-\pi n,\pi-\pi n)$. This is how we read the
positivity condition of \cite[\S4.2]{hkkp26} on the branches of
\cite[(5.15)]{hkkp26}; in \cite[\S5.4]{hkkp26} the phase is written with the
functional calculus, the two readings agree when $\Omega$ is normal, and for
complexes $\Omega$ is not normal in general \cite[Example~2.3]{wue26a}. Two
properties are read off from \eqref{eq:general}:
\begin{enumerate}
\item[(O1)] for an orthogonal direct sum of metrized complexes, $\Omega$ and
 $\Phi$ are the direct sums of those of the summands;
\item[(O2)] if $\mathsf F[-n]$ denotes the complex $\mathsf F$ with all
 degrees raised by $n$, then $\Omega_{\mathsf F[-n]}=(-1)^n\Omega_{\mathsf F}$
 and $\Phi_{\mathsf F[-n]}=\Phi_{\mathsf F}-\pi n$.
\end{enumerate}
The \emph{phase flow} is the flow
\begin{equation}\label{eq:flow}
 h^{-1}\dot h=-\Phi
\end{equation}
of the metrics; it is the minimizing flow of \cite[\S4.1]{hkkp26} for this
phase, and we use a neutral name because the mass, defined below, can
increase along it, already for two-term complexes over $K_m$
\cite[Theorem~1.1]{wue26a}. It exists for all $t\geq0$ from every initial
metric, because the vector field is real-analytic and the velocity
$h^{-1}\dot h$ is bounded \cite[Proposition~2.2]{wue26c}. The \emph{mass} of a metric is
$\mass(h)=\sum_{i,n}\tr|\Omega_i^n|$, where $\Omega_i^n$ is the restriction
of $\Omega$ to $E_i^n$, and a metric is \emph{harmonic} if $\Phi$ is a
scalar. Let each vector $e_k(t)$ of a frame of $E_i^n$ solve
$\dot e_k=\frac12\Phi e_k$, starting from an $h(0)$-orthonormal frame; the
frame is $h(t)$-orthonormal, because $h^{-1}\dot h=-\Phi$ and $\Phi$ is
self-adjoint for $h$. Differentiating the matrices in these frames shows
that the matrix of each structure map $T\in\{d_i,\rho_a\}$ evolves by
\begin{equation}\label{eq:frameflow}
 \dot T=\tfrac12\bigl(T\,\Phi_{\mathrm{source}}-\Phi_{\mathrm{target}}\,T\bigr),
\end{equation}
with the phases on the source and the target of $T$. By (O2) the
trajectories of the flow of $\mathsf F[-n]$ are those of $\mathsf F$
multiplied by $e^{\pi nt}$, and \eqref{eq:frameflow} is the same for both.

For a complex with a surjective differential at $i_0$, \eqref{eq:general}
gives the blocks \eqref{eq:blocks} below \cite[Theorem~6.1(a)]{wue26c}.
Indeed, at $i_0$ one has $\Delta_{i_0}=B\oplus L$ on $V\oplus U$, and
$\Delta_j=0$ for $j\neq i_0$; $\eta_a=-\rho_aD^\dagger$ on $U$ for the arrows
$a$ out of $i_0$, and $\eta_a=0$ for all other arrows; the matrix
$G_{\mathrm{in}}$ has its range in $\ker D$, so that
$K_B(G_{\mathrm{in}})=G_{\mathrm{in}}$; and
$\int e^{-(1-x-y)L}dx\,dy=(L-1+e^{-L})L^{-2}$ over the simplex, so that
$D^\dagger(L-1+e^{-L})L^{-2}D=1-R(B)$. Hence
\begin{equation}\label{eq:blocks}
 \begin{gathered}
 \Omega_V=z_\ast e^{-B}+M_V,\qquad M_V:=K_B(C_{\mathrm{out}})-G_{\mathrm{in}},\qquad
 \Omega_U=J-z_\ast e^{-L},\\
 \Omega_j=z_j+M_j,\qquad
 M_j:=\sum_{a:j\to\cdot}\rho_a^\dagger\rho_a
 -\sum_{\substack{a:k\to j\\k\neq i_0}}\rho_a\rho_a^\dagger
 -\sum_{a:i_0\to j}\rho_aR(B)\rho_a^\dagger\qquad(j\neq i_0).
 \end{gathered}
\end{equation}
The matrices $M_V$, $J$ and $M_j$ are Hermitian and do not depend on the
charges; the $\Omega_j$ are normal; and $\tr\Omega=Z$
\cite[Lemma~5.17]{hkkp26}. The phases are $\Phi_V$ and $\Phi_j$, with
spectrum in $(0,\pi)$, and $\Phi_U$, with spectrum in $(-\pi,0)$; the mass is
$\mass(h)=\tr|\Omega_V|+\tr|\Omega_U|+\sum_{j\neq i_0}\tr|\Omega_j|$.

\subsection{The adapted splitting}
Let $\mathsf E$ have a surjective differential at $i_0$, and fix a metric.
Split
$V=\ker D\oplus(\ker D)^\perp$ orthogonally and write the blocks of $\Phi_V$
as $\Phi_{HH}$, $\Phi_{HA}$, $\Phi_{AH}$, $\Phi_{AA}$, the index $H$ referring
to $\ker D$ and $A$ to its complement. In orthonormal frames adapted to the
splitting,
\[
 D=(0,F),\qquad\rho_a=(H_a,T_a)\ \ (a:i_0\to j),\qquad
 \rho_a=\begin{pmatrix}R_a\\0\end{pmatrix}\ \ (a:k\to i_0),
\]
with $F$ invertible. Put $X_a=T_aF^{-1}$ for the arrows out of $i_0$. The
splitting depends on the metric. As in \cite[\S6.1]{wue26a}, the frame of $V$
is kept adapted by the additional rotation with the anti-Hermitian generator
$\mathsf a=\frac12\bigl(\begin{smallmatrix}0&\Phi_{HA}\\-\Phi_{AH}&0\end{smallmatrix}\bigr)$,
and \eqref{eq:frameflow} becomes
\begin{equation}\label{eq:adapted}
 \begin{gathered}
 \dot F=\tfrac12(F\Phi_{AA}-\Phi_UF),\qquad
 \dot H_a=\tfrac12(H_a\Phi_{HH}-\Phi_jH_a),\\
 \dot X_a=\tfrac12(X_a\Phi_U-\Phi_jX_a)+H_a\Phi_{HA}F^{-1},\qquad
 \dot R_a=\tfrac12(R_a\Phi_k-\Phi_{HH}R_a).
 \end{gathered}
\end{equation}
Indeed, the rotation changes the right sides of \eqref{eq:frameflow} by
$+D\mathsf a$, $+\rho_a\mathsf a$ for the maps out of $V$ and by
$-\mathsf a\rho_a$ for the maps into $V$; the first block of
$\frac12(D\Phi_V-\Phi_UD)+D\mathsf a$ vanishes, the first two equations and
$\dot T_a=\frac12(T_a\Phi_{AA}-\Phi_jT_a)+H_a\Phi_{HA}$ are its second block
and the blocks of $\frac12(\rho_a\Phi_V-\Phi_j\rho_a)+\rho_a\mathsf a$, the
third follows from $X_a=T_aF^{-1}$, and the second block of
$\frac12(\rho_a\Phi_k-\Phi_V\rho_a)-\mathsf a\rho_a$ vanishes.

The arrows of the cohomology representation, for the induced metrics, have
the matrices $\beta_a=H_a$, $R_a$ or $\rho_a$ according as $a$ leaves $i_0$,
enters $i_0$ or does not meet $i_0$. Its \emph{moment map} is
\[
 \mu_i=\sum_{a:i\to\cdot}\beta_a^\dagger\beta_a-\sum_{a:\cdot\to i}\beta_a\beta_a^\dagger
 \qquad(i\in Q_0),
\]
outgoing minus incoming arrows; its endomorphism $\Omega$ is
$\bigoplus_i(z_i+\mu_i)$, which is normal, with phases $\arg(z_i+\mu_i)$.
Put
\[
 q_{\mathrm{out}}=\sum_{a:i_0\to\cdot}\norm{H_a}^2,\qquad
 s=\sum_{a:i_0\to\cdot}\norm{X_a}^2,\qquad
 Y_j=\sum_{a:i_0\to j}X_a(1-e^{-L})X_a^\dagger\geq0 .
\]
Since $B=0\oplus F^\dagger F$ and $FR(F^\dagger F)F^\dagger=1-e^{-L}$, the
blocks \eqref{eq:blocks} are, in the adapted splitting,
\begin{equation}\label{eq:adaptedblocks}
 \Omega_V=\begin{pmatrix}z_\ast+\mu_{i_0}&E_R\\E_R^\dagger&z_\ast e^{-F^\dagger F}+K_{F^\dagger F}(C_T)\end{pmatrix},
 \qquad\Omega_j=z_j+\mu_j-Y_j,
\end{equation}
with $E_R=\sum_aH_a^\dagger T_aR(F^\dagger F)$ and $C_T=\sum_aT_a^\dagger T_a$,
the sums over the arrows out of $i_0$; moreover $C_T=F^\dagger SF$ and
$DC_{\mathrm{out}}D^\dagger=L\,S\,L$ with $S=\sum_aX_a^\dagger X_a$,
$\tr S=s$. Let
\[
 \delta_{i_0}=\Phi_{HH}-\arg(z_\ast+\mu_{i_0}),\qquad
 \delta_j=\Phi_j-\arg(z_j+\mu_j)\quad(j\neq i_0)
\]
be the differences between the phases of the complex on the cohomology and
the phases of the cohomology representation. We say that the exact part
\emph{decays} along a trajectory if $\lambda\to\infty$ and
\[
 \norm{\Omega_V-(z_\ast+\mu_{i_0})\oplus0}_1+\norm{\Omega_U}_1
 +\sum_{j\neq i_0}\norm{\Omega_j-(z_j+\mu_j)}_1\to0,
\]
where $\oplus\,0$ refers to $V=\ker D\oplus(\ker D)^\perp$: the blocks of
$\Omega$ approach those of the cohomology representation with its induced
metrics. If $s=0$, then $E_R=0$, $C_T=0$, $J=0$ and $Y_j=0$: the complex is
the orthogonal direct sum of its cohomology representation and of the
complex $(\ker D)^\perp\to U$ at $i_0$, and all $\delta_i$ vanish.

\subsection{Results used}
We use the following results of \cite{wue26a,wue26b}; they are statements
about matrices.

\begin{proposition}\label{prop:used}
The following inequalities hold.
\begin{enumerate}
\item[(i)] \cite[Theorem~3.1]{wue26b} For Hermitian matrices $L>0$ and
 $M\geq0$, $\tr J_L(M)_-\leq\frac25\tr J_L(M)$.
\item[(ii)] \cite[Corollary~4.2]{wue26b} Let $P>0$ and $K$ be Hermitian
 matrices and $z$ a complex number with $\Ima z>0$, and let $\Psi$ be the
 polar phase of $zP+K$. Then $\tr(K\Psi)\leq(\arg z)\tr K$.
\item[(iii)] \cite[Lemma~3.1]{wue26a} Let
 $T=\bigl(\begin{smallmatrix}A&E\\E^\dagger&C\end{smallmatrix}\bigr)$ have
 positive definite imaginary part, let the least singular value of $A$ be at
 least $b>0$, and let $\Psi$ and $\psi_A$ be the polar phases of $T$ and $A$.
 Then $\norm{\Psi_{12}}\leq2\pi\norm E/b$, and if $\norm E/b\leq\frac18$, then
 $\norm{\Psi_{11}-\psi_A}\leq4\pi\norm E/b$.
\item[(iv)] \cite[Lemma~3.2]{wue26a}, in the form of
 \cite[Proposition~2.3(v)]{wue26b}. Let
 $T_\circ=|T_\circ|\,e^{\ii\Psi_\circ}$ be an invertible matrix in polar form,
 with $\Psi_\circ$ Hermitian, and $T=T_\circ+M$ a matrix with $\Ima T>0$ and
 polar phase $\Psi$. If $v$ is a unit vector with $|T_\circ|v=pv$ and
 $\Psi_\circ v=\psi v$, where $p>0$ and $0<\psi<\pi$, then
 $|\inner v{(\Psi-\psi)v}|\leq\pi\sqrt{\norm M_1/p}$.
\item[(v)] \cite[Lemma~4.2 and \S6.3]{wue26a}, in the form of
 \cite[Proposition~2.3(ii)]{wue26b}. Let $F$ be an invertible $r\times r$
 matrix, $L=FF^\dagger$, and $S\geq0$ a Hermitian $r\times r$ matrix. Then,
 with $\kappa_0=\pi/(2e)$,
 $\norm{K_{F^\dagger F}(F^\dagger SF)}_1\leq r\kappa_0\tr S$ and
 $\norm{J_L(LSL)}_1\leq r\tr S$.
\end{enumerate}
\end{proposition}

We also use two results of \cite{wue26c} on complexes.

\begin{proposition}[\cite{wue26c}]\label{prop:usedc}
\begin{enumerate}
\item[(a)] \cite[Theorem~3.4]{wue26c} Let $\mathsf E$ be a complex with a
 surjective differential at $i_0$, and for a metric $h$ let
 \begin{equation}\label{eq:primitive}
 \cS_{\CC}(h)=\sum_iz_i\log\det h_i-z_\ast\log\det h_U-\sum_a\norm{\rho_a}^2
 +\int_0^1\bigl(\tr\Omega_U^{(x)}-\tr\Omega_U^{(0)}\bigr)\frac{dx}x,
 \end{equation}
 where $h_{i_0}$, $h_U$ and $h_i$ are the metrics of $V$, $U$ and $W_i$, the
 determinants are Gram determinants in fixed bases, and $\Omega^{(x)}$
 denotes $\Omega$ for the complex with $D$ replaced by $\sqrt x\,D$. Then
 $\cS_{\CC}$ is a complex-valued primitive of $\Omega$:
 $d\cS_{\CC}(\xi)=\tr(\Omega_V\xi_V)+\tr(\Omega_U\xi_U)+\sum_j\tr(\Omega_j\xi_j)$
 for $\xi=h^{-1}\delta h$.
\item[(b)] \cite[Theorem~5.1]{wue26c} Let $\mathsf E^\vee$ be the dual
 complex over the opposite quiver: the dual spaces in the opposite degrees,
 the transposed maps, the dual metrics, and the charges $-\overline{z_i}$.
 Then $h(t)\mapsto e^{-\pi t}(h(t)^{-1})^{\mathsf T}$ maps trajectories of
 the phase flow of $\mathsf E$ to trajectories of the phase flow of
 $\mathsf E^\vee$; in dual orthonormal frames $\Omega$ is replaced by
 $-\overline\Omega$ and $\Phi$ by $\pi-\overline\Phi$; and the mass is
 preserved.
\end{enumerate}
\end{proposition}

\section{An a priori bound for the arrows}\label{sec:energy}

\begin{definition}\label{def:energy}
For an arrow $a$ put
\[
 e_a=\norm{\rho_a}^2\quad\text{if }\mathsf s(a)\neq i_0,\qquad
 e_a=\tr\bigl(\rho_aR(B)\rho_a^\dagger\bigr)\quad\text{if }\mathsf s(a)=i_0 .
\]
The \emph{arrow energy} of the metrized complex is $\cN=\sum_{a\in Q_1}e_a$.
\end{definition}

For an arrow $a$ out of $i_0$ we split $e_a=e_a^K+e_a^J$ with
\[
 e_a^K=\tr\bigl(\rho_a^\dagger\rho_a\,e^{-B}\bigr)\geq0,\qquad
 e_a^J=\tr\bigl(\rho_a^\dagger\rho_a\,(R(B)-e^{-B})\bigr)\geq0 .
\]
Then
\begin{equation}\label{eq:traces}
 \tr K_B(C_{\mathrm{out}})=\sum_{a:i_0\to\cdot}e_a^K,\qquad
 \tr J=\sum_{a:i_0\to\cdot}e_a^J,\qquad
 \cN=\sum_a\norm{\beta_a}^2+\sum_{j\neq i_0}\tr Y_j .
\end{equation}
The first identity holds because $\tr K_B(G)=\tr(Ge^{-B})$. For the second,
let $\omega(\ell)=\int_0^1we^{-w\ell}dw=(1-e^{-\ell}-\ell e^{-\ell})/\ell^2$;
then
\[
 \tr J=\tr\bigl(DC_{\mathrm{out}}D^\dagger\,\omega(L)\bigr),\qquad
 D^\dagger\omega(L)D=R(B)-e^{-B}.
\]
The third follows from
$\rho_aR(B)\rho_a^\dagger=H_aH_a^\dagger+X_a(1-e^{-L})X_a^\dagger$. In
particular
\begin{equation}\label{eq:energybounds}
 \sum_a\norm{\beta_a}^2\leq\cN,\qquad\tr J\leq\cN,\qquad
 (1-e^{-\lambda})\,s\leq\sum_j\tr Y_j\leq\cN,\qquad\norm{\mu_i}_1\leq\cN .
\end{equation}

\begin{lemma}\label{lem:energyderivative}
For all charges, along the phase flow,
\[
 \frac{d\cN}{dt}=\tr(M_V\Phi_V)+\tr(J\Phi_U)+\sum_{j\neq i_0}\tr(M_j\Phi_j).
\]
\end{lemma}

\begin{proof}
In the function \eqref{eq:primitive} of Proposition~\ref{prop:usedc}(a), by
\eqref{eq:blocks},
\[
 \tr\Omega_U^{(x)}=\tr J_{xL}(xDC_{\mathrm{out}}D^\dagger)-z_\ast\tr e^{-xL},\qquad
 \tr J_{xL}(xDC_{\mathrm{out}}D^\dagger)=x\tr\bigl(DC_{\mathrm{out}}D^\dagger\omega(xL)\bigr).
\]
Since $\int_0^1\omega(x\ell)\,dx=(1-R(\ell))/\ell$ and
$D^\dagger\frac{1-R(L)}LD=1-R(B)$, the part of $\cS_{\CC}$ that does not
contain the charges is
\[
 -\sum_a\norm{\rho_a}^2+\tr C_{\mathrm{out}}-\tr\bigl(C_{\mathrm{out}}R(B)\bigr)=-\cN .
\]
Both sides of the identity for $d\cS_{\CC}$ are affine functions of the
charges, so the terms without charges agree:
$d\cN(\xi)=-\tr(M_V\xi_V)-\tr(J\xi_U)-\sum_j\tr(M_j\xi_j)$. Put $\xi=-\Phi$.
\end{proof}

We use the following elementary inequality: for a complex number $\zeta$ in
the upper half-plane and a real number $u\geq0$,
\begin{equation}\label{eq:scalar}
 u\bigl(\pi-\arg(\zeta-u)\bigr)\leq\max\{\pi\max(\Rea\zeta,0),\,2\Ima\zeta\}.
\end{equation}
It is trivial for $u=0$; let $u>0$. If $\Rea\zeta\leq u/2$ then
$\pi-\arg(\zeta-u)=\arctan\frac{\Ima\zeta}{u-\Rea\zeta}\leq\frac{2\Ima\zeta}u$,
and otherwise $u<2\Rea\zeta$ and $\pi-\arg(\zeta-u)\leq\frac\pi2$ for
$u\geq\Rea\zeta$, while for $u<\Rea\zeta$ the left side is at most
$\pi u<\pi\Rea\zeta$.

\begin{lemma}\label{lem:three}
For every metric and all charges:
\begin{enumerate}
\item[(a)] $\tr(M_V\Phi_V)\leq\alpha_\ast\tr M_V
 =\alpha_\ast\bigl(\sum_{a:i_0\to\cdot}e_a^K-\sum_{a:\cdot\to i_0}e_a\bigr)$;
\item[(b)] $\tr(J\Phi_U)\leq\pi\tr J_-\leq\frac{2\pi}5\tr J
 =\frac{2\pi}5\sum_{a:i_0\to\cdot}e_a^J$, and $\tr(J\Phi_U)\leq0$ if $r\leq1$;
\item[(c)] for $j\neq i_0$, with the number $b_j:=\max\{\pi|x_j|,\,2y_j\}$,
 \[
 \tr(M_j\Phi_j)\leq n_jb_j-\pi\tr(M_j)_-,\qquad
 \tr(M_j)_-\geq\Bigl(\sum_{a:\cdot\to j}e_a-\sum_{a:j\to\cdot}e_a\Bigr)_+ .
 \]
\end{enumerate}
\end{lemma}

\begin{proof}
(a) Proposition~\ref{prop:used}(ii) with $zP=z_\ast e^{-B}$ and $K=M_V$; and
$\tr G_{\mathrm{in}}=\sum_{a:\cdot\to i_0}e_a$.

(b) Since $-\pi<\Phi_U<0$, $\tr(J_+\Phi_U)\leq0$ and
$-\tr(J_-\Phi_U)\leq\pi\tr J_-$; then apply
Proposition~\ref{prop:used}(i) to $M=DC_{\mathrm{out}}D^\dagger\geq0$. If
$r=1$, $J$ is a nonnegative number.

(c) The matrix $\Omega_j=z_j+M_j$ is normal and $\Phi_j=g_j(M_j)$ with
$g_j(u)=\arg(z_j+u)\in(0,\pi)$ for real $u$. For $u\geq0$,
$\arg(z_j+u)=\pi-\arg(-\overline{z_j}-u)$, so \eqref{eq:scalar} with
$\zeta=-\overline{z_j}$ gives $u\,g_j(u)\leq b_j$. For $u<0$,
$u\,g_j(u)=-\pi|u|+|u|\bigl(\pi-\arg(z_j-|u|)\bigr)\leq-\pi|u|+b_j$ by
\eqref{eq:scalar} with $\zeta=z_j$. Summing over the eigenvalues of $M_j$ gives
the first inequality. Finally $\tr(M_j)_-\geq(-\tr M_j)_+$, and
$-\tr M_j=\sum_{a:\cdot\to j}e_a-\sum_{a:j\to\cdot}e_a$ by \eqref{eq:blocks}.
\end{proof}

\subsection{Height functions}
The estimates of Lemma~\ref{lem:three} at the different vertices are combined
with coefficients that increase along the arrows. Let
\[
 \alpha_\ast^+=\max\{\alpha_\ast,\tfrac{2\pi}5\}\ \text{ if }r\geq2,\qquad
 \alpha_\ast^+=\alpha_\ast\ \text{ if }r\leq1 .
\]
A \emph{height function} is a function $w:Q_0\setminus\{i_0\}\to[0,1]$ such
that
\[
 w_{\mathsf t(a)}>w_{\mathsf s(a)}\ \ \text{if }\mathsf s(a),\mathsf t(a)\neq i_0,\qquad
 \pi w_{\mathsf t(a)}>\alpha_\ast^+\ \ \text{if }\mathsf s(a)=i_0,\qquad
 \pi w_{\mathsf s(a)}<\alpha_\ast\ \ \text{if }\mathsf t(a)=i_0 .
\]
Height functions exist. Let $\mathcal A$ be the set of vertices different from $i_0$
with a path to $i_0$, $\mathcal D$ the set of vertices different from $i_0$
reached by a path from $i_0$, and $\mathcal R$ the set of the remaining
vertices different from $i_0$. Since $Q$ has no oriented cycle,
$\mathcal A\cap\mathcal D=\emptyset$, there is no arrow from
$\mathcal R\cup\mathcal D$ to $\mathcal A$, and none from $\mathcal D$ to
$\mathcal R$. So there is a linear order of $Q_0\setminus\{i_0\}$, compatible
with the arrows, in which the vertices of $\mathcal A$ come first and those
of $\mathcal D$ last; let $w$ be increasing for this order, with values in
$(0,\frac{\alpha_\ast}{2\pi})$ on $\mathcal A$, in
$(\frac{\alpha_\ast}{2\pi},\frac{\pi+\alpha_\ast^+}{2\pi})$ on $\mathcal R$ and in
$(\frac{\pi+\alpha_\ast^+}{2\pi},1)$ on $\mathcal D$. For a height function
$w$ put
\begin{equation}\label{eq:rhoQ}
 \begin{gathered}
 \gamma_w=\min\Bigl\{\min_{\substack{a\in Q_1\\\mathsf s(a),\mathsf t(a)\neq i_0}}\pi\bigl(w_{\mathsf t(a)}-w_{\mathsf s(a)}\bigr),\
 \min_{a:i_0\to\cdot}\bigl(\pi w_{\mathsf t(a)}-\alpha_\ast^+\bigr),\
 \min_{a:\cdot\to i_0}\bigl(\alpha_\ast-\pi w_{\mathsf s(a)}\bigr)\Bigr\}>0,\\
 c_Q=\sum_{j\neq i_0}n_jb_j,
 \end{gathered}
\end{equation}
where a minimum over an empty set is omitted; if $Q$ has no arrows, put
$\gamma_w=1$.

\begin{theorem}[A priori bound]\label{thm:energy}
Let $\mathsf E$ be a complex with a surjective differential at $i_0$ and
$w$ a height function. For all charges in the upper half-plane and every
trajectory of the phase flow,
\[
 \frac{d\cN}{dt}\leq-\gamma_w\,\cN+c_Q,\qquad
 \cN(t)\leq\cN(0)e^{-\gamma_wt}+\frac{c_Q}{\gamma_w}\bigl(1-e^{-\gamma_wt}\bigr).
\]
Hence $\cN(t)\leq\max\{\cN(0),c_Q/\gamma_w\}$ for all $t\geq0$, and if $c_Q>0$,
\[
 \cN(t)\leq\bar\cN:=\frac{2c_Q}{\gamma_w}\qquad\text{for }t\geq
 t_0:=\frac1{\gamma_w}\log^+\frac{\gamma_w\cN(0)}{c_Q},
\]
where $\log^+x=\max(\log x,0)$ for $x>0$ and $\log^+0=0$.
\end{theorem}

\begin{proof}
For $w_j\in[0,1]$ one has $\tr(M_j)_-\geq w_j\,(-\tr M_j)$, by
Lemma~\ref{lem:three}(c). So Lemmas~\ref{lem:energyderivative}
and~\ref{lem:three} give
\[
 \frac{d\cN}{dt}\leq c_Q-\pi\sum_{j\neq i_0}w_j\Bigl(\sum_{a:\cdot\to j}e_a-\sum_{a:j\to\cdot}e_a\Bigr)
 +\alpha_\ast\!\!\sum_{a:i_0\to\cdot}\!\!e_a^K+\tfrac{2\pi}5\!\!\sum_{a:i_0\to\cdot}\!\!e_a^J
 -\alpha_\ast\!\!\sum_{a:\cdot\to i_0}\!\!e_a,
\]
where for $r\leq1$ the term with $\frac{2\pi}5$ is absent. Collect the terms
arrow by arrow. An arrow $a$ with $\mathsf s(a),\mathsf t(a)\neq i_0$ has the coefficient
$-\pi(w_{\mathsf t(a)}-w_{\mathsf s(a)})$. An arrow $a$ out of $i_0$ has $-\pi w_{\mathsf t(a)}e_a$
from the first sum and at most $\alpha_\ast^+(e_a^K+e_a^J)=\alpha_\ast^+ e_a$ from
the next two. An arrow $a$ into $i_0$ has $+\pi w_{\mathsf s(a)}e_a-\alpha_\ast e_a$.
Hence $\frac{d\cN}{dt}\leq c_Q-\gamma_w\sum_ae_a$. The second inequality
follows by integration, and
$\cN(0)e^{-\gamma_wt}\leq c_Q/\gamma_w$ for $t\geq t_0$.
\end{proof}

\begin{remark}\label{rem:energy}
(1) By \eqref{eq:energybounds} the theorem bounds the arrows of the
cohomology representation, its moment map, $\tr J$, and
$s\leq\cN/(1-e^{-\lambda})$. The bound $\bar\cN$ depends only on the charges,
the dimensions $n_j$ and the height function; the initial metric enters only
through
$t_0$. The height function constructed above is determined by the quiver,
the vertex $i_0$ and the numbers $\alpha_\ast$, $\alpha_\ast^+$ once the
linear order and the values are fixed by a rule, for instance values equally
spaced in each of the three intervals; with such a rule the constants of
Theorem~\ref{thm:energy} and of \eqref{eq:constants} depend only on the
quiver, $i_0$, the charges and the dimensions, except $\kappa_1$ and the
times $t_0$, $t_1$, $t_2$, which depend in addition on the initial metric.

(2) For $Q=K_m$ and $i_0=0$ the only condition on $w$ is
$\pi w_1>\alpha_\ast^+$, and $w_1=1$ gives $\gamma_w=\pi-\alpha_\ast^+$, the
constant of \cite[Theorem~5.4]{wue26b}. For a vertex $j$ that is a sink the
number $b_j$ may be replaced by $\max\{\pi\max(x_j,0),2y_j\}$, and for a source
by $\max\{\pi\max(-x_j,0),2y_j\}$, since only one sign of the eigenvalues of
$M_j$ occurs.

(3) For a quiver with an oriented cycle there is no such bound in general: a
loop acting by a multiple $x$ of a unitary matrix has zero moment map for
every $x$, the phases are then constant, and the arrow does not change along
the flow.
\end{remark}

For a representation no vertex is distinguished, and the bound takes the
following form. It is used in Theorem~\ref{thm:representations}.

\begin{proposition}[Representations]\label{prop:reparrows}
Let $\mathsf R$ be a nonzero representation of a finite quiver $Q$ without
oriented cycles, with dimension vector $(n_i)$ and any metric. Let $\ell_Q\geq1$ be
the number of arrows of a longest path in $Q$; if $Q$ has no arrows, put
$\ell_Q=1$. For all charges in the upper half-plane, the sum
$\cN=\sum_a\norm{\rho_a}^2$ of the squared norms of the arrows satisfies,
along the phase flow,
\[
 \frac{d\cN}{dt}\leq-\frac\pi{\ell_Q}\,\cN+c_{\mathsf R},\qquad
 c_{\mathsf R}=\sum_in_i\max\{\pi|x_i|,2y_i\}.
\]
Hence $\cN(t)\leq\bar\cN_{\mathsf R}:=2\ell_Qc_{\mathsf R}/\pi$ for
$t\geq t_{\mathsf R}:=\frac{\ell_Q}\pi\log^+\frac{\pi\cN(0)}{\ell_Qc_{\mathsf R}}$.
\end{proposition}

\begin{proof}
By \eqref{eq:frameflow},
$\frac d{dt}\norm{\rho_a}^2=\tr(\rho_a^\dagger\rho_a\Phi_{\mathsf s(a)})
-\tr(\rho_a\rho_a^\dagger\Phi_{\mathsf t(a)})$, so
$\frac{d\cN}{dt}=\sum_i\tr(\mu_i\Phi_i)$ with the moment map
$\mu_i=\sum_{a:i\to\cdot}\rho_a^\dagger\rho_a-\sum_{a:\cdot\to i}\rho_a\rho_a^\dagger$.
The block $\Omega_i=z_i+\mu_i$ is normal and $\Phi_i=\arg(z_i+\mu_i)$, so the
argument of Lemma~\ref{lem:three}(c) gives, for every vertex $i$ and every
$w_i\in[0,1]$,
\[
 \tr(\mu_i\Phi_i)\leq n_i\max\{\pi|x_i|,2y_i\}-\pi\tr(\mu_i)_-
 \leq n_i\max\{\pi|x_i|,2y_i\}
 -\pi w_i\Bigl(\sum_{a:\cdot\to i}\norm{\rho_a}^2-\sum_{a:i\to\cdot}\norm{\rho_a}^2\Bigr).
\]
Let $\ell(i)$ be the number of arrows of a longest path ending at $i$, and
$w_i=\ell(i)/\ell_Q$; then $w_{\mathsf t(a)}-w_{\mathsf s(a)}\geq1/\ell_Q$ for
every arrow $a$. Summing over $i$ gives
$\frac{d\cN}{dt}\leq c_{\mathsf R}-\pi\sum_a(w_{\mathsf t(a)}-w_{\mathsf s(a)})\norm{\rho_a}^2
\leq c_{\mathsf R}-\frac\pi{\ell_Q}\cN$, and the last statement follows by
integration.
\end{proof}

\section{Growth of the differential and decay of the exact part}\label{sec:decay}

In this section $\mathsf E$ is a complex with a surjective differential at
$i_0$ and $r\geq1$. The estimates of Lemma~\ref{lem:estimates} and the proof
of Theorem~\ref{thm:decay} follow those of \cite{wue26b} for $K_m$; what is
new is that the block $z_\ast+\mu_{i_0}$ is no longer $z_\ast$ plus a
positive matrix, and that the phases at the targets of the arrows out of
$i_0$ are bounded below only through the a priori bound. Recall $\lambda=\lambda_{\min}(L)=\lambda_{\min}(F^\dagger F)$, and
let
\[
 \varpi_\star=\min\bigl\{\lambda_{\min}(\Phi_j):\ j=\mathsf t(a)\ \text{for an arrow }a\text{ out of }i_0\bigr\};
\]
only targets with $W_j\neq0$ enter this minimum. If there are none,
then $s=0$, Lemma~\ref{lem:estimates}(iii) reads
$0\leq0$, and $\varpi_\star$, $\bar\varpi$ and the constants
\eqref{eq:constants} are not used.
Let $\Eone(b)=\int_1^\infty e^{-xb}\frac{dx}x$ be the exponential integral;
for $b>0$,
\begin{equation}\label{eq:E1}
 \frac{e^{-b}}{1+b}\leq\Eone(b)\leq e^{-b}\log\bigl(1+\tfrac1b\bigr),\qquad
 \Eone'(b)=-\frac{e^{-b}}b ,
\end{equation}
because the differences of the three functions tend to $0$ at infinity and
have derivatives of constant sign.

\begin{lemma}\label{lem:estimates}
For all charges and every metric:
\begin{enumerate}
\item[(i)] $\lambda$ is nondecreasing along the flow; its right derivative is
 the minimum of
 $\lambda\bigl(\inner v{\Phi_{AA}v}-\inner u{\Phi_Uu}\bigr)$ over the unit
 vectors $v$ with $F^\dagger Fv=\lambda v$, where $u=Fv/\sqrt\lambda$.
\item[(ii)] $\norm{E_R}\leq\sqrt{q_{\mathrm{out}}s/\lambda}$ and
 $\sum_j\norm{Y_j}_1\leq s$; hence, with $\mathsf d=\min(n_{i_0},r)$ and
 $\mathcal E_V:=\Omega_V-\diag(z_\ast+\mu_{i_0},\ z_\ast e^{-F^\dagger F})$,
 \[
 \norm{\mathcal E_V}_1\leq2\sqrt{\mathsf d\,q_{\mathrm{out}}s/\lambda}+r\kappa_0s .
 \]
\item[(iii)] $\dot s\leq\bigl(-\varpi_\star+\frac{4\pi q_{\mathrm{out}}}{y_\ast\lambda}\bigr)s$
 along the flow. Moreover $\lambda_{\min}(\Phi_j)\geq\arg(z_j+\cN)$ for
 each $j\neq i_0$ with $W_j\neq 0$. At every sink with $W_j\neq 0$ one has $\lambda_{\min}(\Phi_j)\geq\alpha_j$.
\item[(iv)] If $\sqrt{q_{\mathrm{out}}s/\lambda}\leq y_\ast/8$, then
 $\norm{\delta_{i_0}}\leq\frac{4\pi}{y_\ast}\sqrt{q_{\mathrm{out}}s/\lambda}$;
 and $\norm{\delta_j}\leq\norm{Y_j}/y_j\leq s/y_j$ for $j\neq i_0$.
\item[(v)] Along the flow
 \[
 \frac d{dt}\tr\Eone(L)\leq-\tr\bigl(e^{-L}|\Phi_U|\bigr),\qquad
 \tr\bigl(e^{-L}|\Phi_U|\bigr)\geq
 \frac{y_\ast(\tr e^{-L})^2}{(y_\ast+|x_\ast|)\tr e^{-L}+\tr J_+} .
 \]
\end{enumerate}
\end{lemma}

\begin{proof}
(i) By \eqref{eq:adapted},
$\frac d{dt}(F^\dagger F)=\frac12(\Phi_{AA}F^\dagger F+F^\dagger F\Phi_{AA})-F^\dagger\Phi_UF$.
For a unit vector $v$ with $F^\dagger Fv=\lambda v$ this gives
$\inner v{\frac d{dt}(F^\dagger F)v}=\lambda\inner v{\Phi_{AA}v}-\inner{Fv}{\Phi_UFv}$,
which is positive because $\Phi_{AA}>0>\Phi_U$; and the right derivative of
the least eigenvalue of a differentiable family of Hermitian matrices is the
minimum of these numbers over the unit vectors of the eigenspace.

(ii) Since $E_R=\sum_aH_a^\dagger X_a\,FR(F^\dagger F)$ and
$\norm{FR(F^\dagger F)}_{\mathrm{op}}\leq\lambda^{-1/2}$, the Cauchy--Schwarz
inequality gives the first bound. $\sum_j\tr Y_j=\tr\bigl((1-e^{-L})S\bigr)\leq s$.
The off-diagonal part of $\mathcal E_V$ has rank at most $2\mathsf d$ and
trace norm $2\norm{E_R}_1\leq2\sqrt{\mathsf d}\norm{E_R}$, and its remaining
block is $K_{F^\dagger F}(C_T)=K_{F^\dagger F}(F^\dagger SF)$; use
Proposition~\ref{prop:used}(v).

(iii) By \eqref{eq:adapted},
\[
 \dot s=\sum_a\tr(X_a^\dagger X_a\Phi_U)-\sum_a\tr(X_a^\dagger\Phi_{\mathsf t(a)}X_a)
 +2\Rea\sum_a\tr\bigl(X_a^\dagger H_a\Phi_{HA}F^{-1}\bigr).
\]
The first sum is nonpositive, the second is at least $\varpi_\star s$, and the last
term is at most $2\sqrt{s\,q_{\mathrm{out}}/\lambda}\,\norm{\Phi_{HA}}$. In
\eqref{eq:adaptedblocks} the block $z_\ast+\mu_{i_0}$ is normal with
imaginary part $y_\ast$, so its singular values are at least $y_\ast$, and
Proposition~\ref{prop:used}(iii) gives
$\norm{\Phi_{HA}}\leq\frac{2\pi}{y_\ast}\norm{E_R}\leq\frac{2\pi}{y_\ast}\sqrt{q_{\mathrm{out}}s/\lambda}$.
For the phases: $\Phi_j=g_j(M_j)$ with $g_j(u)=\arg(z_j+u)$ decreasing, and
$M_j\leq\sum_{a:j\to\cdot}\rho_a^\dagger\rho_a\leq\cN$ by \eqref{eq:blocks};
for a sink $M_j\leq0$.

(iv) The first statement is Proposition~\ref{prop:used}(iii) for
\eqref{eq:adaptedblocks}. For $j\neq i_0$, $\Phi_j=g_j(\mu_j-Y_j)$ and
$\arg(z_j+\mu_j)=g_j(\mu_j)$ by \eqref{eq:adaptedblocks}, $|g_j'|\leq1/y_j$, and a
function with Lipschitz constant $c$ satisfies
$\norm{g(A)-g(A')}\leq c\norm{A-A'}$ for Hermitian matrices $A,A'$ in the
Hilbert--Schmidt norm, as one sees by expanding in eigenbases of $A$ and
$A'$.

(v) By \eqref{eq:adapted},
$\dot L=F\Phi_{AA}F^\dagger-\frac12(\Phi_UL+L\Phi_U)$, so by \eqref{eq:E1}
\[
 \frac d{dt}\tr\Eone(L)=-\tr\bigl(e^{-F^\dagger F}\Phi_{AA}\bigr)+\tr\bigl(e^{-L}\Phi_U\bigr),
\]
and both terms are nonpositive. Let $(u_k)$ be an orthonormal eigenbasis of
$\Phi_U$ with eigenvalues $\psi_k\in(-\pi,0)$, and put
$e_k=\inner{u_k}{e^{-L}u_k}>0$, $j_k=\inner{u_k}{Ju_k}$. Since
$-\Omega_U=z_\ast e^{-L}-J$ has the polar phase $\Phi_U+\pi$ and a positive
definite modulus, $\inner{u_k}{(z_\ast e^{-L}-J)u_k}$ is a positive multiple
of $e^{\ii(\psi_k+\pi)}$; hence
$\cot(\psi_k+\pi)=(x_\ast e_k-j_k)/(y_\ast e_k)$ and
\[
 |\psi_k|=\operatorname{arccot}\frac{j_k-x_\ast e_k}{y_\ast e_k}
 \geq\operatorname{arccot}\frac{(j_k)_++|x_\ast|e_k}{y_\ast e_k}
 \geq\frac{y_\ast e_k}{(y_\ast+|x_\ast|)e_k+(j_k)_+},
\]
because $\operatorname{arccot}$ is decreasing and
$\operatorname{arccot}x\geq\frac1{1+x}$ for $x\geq0$. Therefore
\[
 \tr\bigl(e^{-L}|\Phi_U|\bigr)=\sum_k|\psi_k|e_k
 \geq\sum_k\frac{y_\ast e_k^2}{(y_\ast+|x_\ast|)e_k+(j_k)_+}
 \geq\frac{y_\ast(\sum_ke_k)^2}{(y_\ast+|x_\ast|)\sum_ke_k+\sum_k(j_k)_+}
\]
by the Cauchy--Schwarz inequality, and $\sum_ke_k=\tr e^{-L}$,
$\sum_k(j_k)_+\leq\tr J_+$.
\end{proof}

\subsection{Constants}
Fix a height function $w$ and let $\gamma_w$, $c_Q$, $\bar\cN$, $t_0$ be as in
Theorem~\ref{thm:energy}. We assume $c_Q>0$, that is, $W_j\neq0$ for some
$j\neq i_0$; otherwise all arrows vanish, Proposition~\ref{prop:sink} below
applies, and we put $\bar\cN=0$ and $t_0=0$. Let $\bar\varpi$ be the least of the
numbers
\[
 \arg(z_j+\bar\cN)\ \ (j\text{ not a sink}),\qquad\alpha_j\ \ (j\text{ a sink}),
\]
over the nonzero target spaces $W_j$ of arrows out of $i_0$; by
Lemma~\ref{lem:estimates}(iii) and Theorem~\ref{thm:energy},
$\varpi_\star\geq\bar\varpi$ for $t\geq t_0$. Let $\lambda_0=\lambda(0)$,
$c_\lambda=1/\log(1+1/\lambda_0)$, $\Xi_0=\tr\Eone(L)|_{t=0}$, and
\begin{equation}\label{eq:constants}
 \begin{gathered}
 \kappa_1=\frac{y_\ast\,c_\lambda^2}{\sqrt2\,|z_\ast|\,c_\lambda\,\Xi_0+\frac75\bar\cN},\qquad
 \Lambda_1=\max\Bigl\{1,\frac{8\pi\bar\cN}{\bar\varpi\,y_\ast}\Bigr\},\\
 t_1=t_0+\frac{(1+\Lambda_1)e^{\Lambda_1}}{\kappa_1},\qquad
 s_1=\frac{\bar\cN}{1-e^{-1}},\qquad
 \mathfrak c_1=2\sqrt{\mathsf d\,\bar\cN s_1}+r\kappa_0s_1,\\
 \mathfrak c_2=\frac\pi{\sqrt{|z_\ast|}}\bigl(\sqrt{\mathfrak c_1}+\sqrt{rs_1}\bigr),\qquad
 \mathfrak c_3=\max\Bigl\{0,2\log\frac{2\mathfrak c_2}\pi\Bigr\},\\
 C_1=\mathfrak c_1+(r+1)s_1+2|z_\ast|re^{\mathfrak c_3},\qquad
 t_2=t_1+\frac2{\bar\varpi}\log^+\frac{8\bar\varpi\,s_1}{\pi y_\ast}.
 \end{gathered}
\end{equation}
In these constants the subscripts are ordinal labels and do not refer to
vertices.

\begin{theorem}[Decay of the exact part]\label{thm:decay}
Let $\mathsf E$ be a complex with a surjective differential at $i_0$ and
$r\geq1$, with at least one arrow from $i_0$ to a vertex $j$ with
$W_j\neq0$; fix a height function $w$ and let the constants be as in
\eqref{eq:constants}. For all
charges in the upper half-plane and along every trajectory of the phase flow:
\begin{enumerate}
\item[(i)] $\lambda(t)\to\infty$; for $t\geq t_1$, $\lambda(t)\geq
 \max\{\Lambda_1,\ \frac{\bar\varpi}4(t-t_1)-\mathfrak c_3\}$ and
 $s(t)\leq s_1e^{-\frac{\bar\varpi}2(t-t_1)}$;
\item[(ii)] for $t\geq t_1$,
 \[
 \norm{\Omega_V-(z_\ast+\mu_{i_0})\oplus0}_1+\norm{\Omega_U}_1
 +\sum_{j\neq i_0}\norm{\Omega_j-(z_j+\mu_j)}_1
 \leq C_1\,e^{-\frac{\bar\varpi}4(t-t_1)};
 \]
 the difference between the mass and the mass $\sum_i\norm{z_i+\mu_i}_1$ of
 the cohomology representation with its induced metrics satisfies the same
 bound;
\item[(iii)] for $t\geq t_2$,
 $\norm{\delta_{i_0}}^2\leq\frac{2\pi\bar\varpi}{y_\ast}\,s$ and
 $\norm{\delta_j}\leq s/y_j$ for $j\neq i_0$.
\end{enumerate}
In particular the exact part decays.
\end{theorem}

The theorem makes no assumption on the cohomology representation.

\begin{proof}
\emph{Growth of $\lambda$.} Put $\Xi=\tr\Eone(L)$. By
Lemma~\ref{lem:estimates}(i) all eigenvalues of $L$ are at least $\lambda_0$
for $t\geq0$, so \eqref{eq:E1} gives $\tr e^{-L}\geq c_\lambda\Xi$. For
$t\geq t_0$, $\tr J_+=\tr J+\tr J_-\leq\frac75\tr J\leq\frac75\bar\cN$ by
Proposition~\ref{prop:used}(i), \eqref{eq:energybounds} and
Theorem~\ref{thm:energy}, and $y_\ast+|x_\ast|\leq\sqrt2|z_\ast|$. The
function $x\mapsto x^2/(cx+d)$ is increasing for $x>0$, and $\Xi$ is
nonincreasing; so Lemma~\ref{lem:estimates}(v) gives, for $t\geq t_0$,
\[
 \dot\Xi\leq-\frac{y_\ast c_\lambda^2\,\Xi^2}{\sqrt2|z_\ast|c_\lambda\,\Xi+\frac75\bar\cN}
 \leq-\kappa_1\Xi^2,
\]
hence $\Eone(\lambda(t))\leq\Xi(t)<1/(\kappa_1(t-t_0))$ for $t>t_0$. For
$t\geq t_1$ this is at most $e^{-\Lambda_1}/(1+\Lambda_1)\leq\Eone(\Lambda_1)$
by \eqref{eq:E1}, so $\lambda(t)\geq\Lambda_1$.

\emph{Decay of $s$.} For $t\geq t_0$ we have $q_{\mathrm{out}}\leq\bar\cN$ and
$\varpi_\star\geq\bar\varpi$, so for $t\geq t_1$,
$\frac{4\pi q_{\mathrm{out}}}{y_\ast\lambda}\leq\frac{\bar\varpi}2$ and
$\dot s\leq-\frac{\bar\varpi}2s$ by Lemma~\ref{lem:estimates}(iii); at $t_1$,
$\lambda\geq1$ and $s\leq\bar\cN/(1-e^{-\lambda})\leq s_1$ by
\eqref{eq:energybounds}.

\emph{Linear growth of $\lambda$.} By Lemma~\ref{lem:estimates}(ii) and
Proposition~\ref{prop:used}(v), for $t\geq t_1$,
$\norm{\mathcal E_V}_1\leq\mathfrak c_1e^{-\frac{\bar\varpi}4(t-t_1)}$ and
$\norm J_1\leq rs_1e^{-\frac{\bar\varpi}2(t-t_1)}$. Let $v$ be a unit vector of
$(\ker D)^\perp$ with $F^\dagger Fv=\lambda v$ and $u=Fv/\sqrt\lambda$. Apply
Proposition~\ref{prop:used}(iv) to $T=\Omega_V$ with
$T_\circ=\diag(z_\ast+\mu_{i_0},\ z_\ast e^{-F^\dagger F})$, whose first
block is normal and invertible and for which
$|T_\circ|v=|z_\ast|e^{-\lambda}v$ and $\Psi_\circ v=\alpha_\ast v$, and to
$T=-\Omega_U=z_\ast e^{-L}-J$ with $T_\circ=z_\ast e^{-L}$ and the vector
$u$:
\[
 |\inner v{(\Phi_V-\alpha_\ast)v}|\leq\pi\sqrt{\frac{\norm{\mathcal E_V}_1e^\lambda}{|z_\ast|}},
 \qquad
 |\inner u{(\Phi_U+\pi-\alpha_\ast)u}|\leq\pi\sqrt{\frac{\norm J_1e^\lambda}{|z_\ast|}} .
\]
By Lemma~\ref{lem:estimates}(i) the right derivative of $\lambda$ satisfies
$\dot\lambda\geq\lambda\bigl[\pi-\mathfrak c_2e^{\lambda/2-\frac{\bar\varpi}8(t-t_1)}\bigr]$
for $t\geq t_1$. Where $\lambda$ is at most
$\frac{\bar\varpi}4(t-t_1)-\mathfrak c_3$, the bracket is at least $\frac\pi2$, so
$\dot\lambda\geq\frac\pi2\lambda\geq\frac\pi2>\frac{\bar\varpi}4$; since
$\lambda(t_1)\geq1>-\mathfrak c_3$, the graph of $\lambda$ cannot cross below
this line. This proves (i).

(ii) By \eqref{eq:adaptedblocks},
$\norm{\Omega_V-(z_\ast+\mu_{i_0})\oplus0}_1\leq\norm{\mathcal E_V}_1+|z_\ast|re^{-\lambda}$,
$\norm{\Omega_U}_1\leq rs+|z_\ast|re^{-\lambda}$ and
$\sum_j\norm{\Omega_j-(z_j+\mu_j)}_1=\sum_j\tr Y_j\leq s$; by (i),
$e^{-\lambda}\leq e^{\mathfrak c_3}e^{-\frac{\bar\varpi}4(t-t_1)}$. The statement
on the masses follows from the triangle inequality for the trace norm.

(iii) For $t\geq t_2$, $s\leq\frac{\pi y_\ast}{8\bar\varpi}$ by (i). Then
$\frac{q_{\mathrm{out}}s}{\lambda y_\ast^2}\leq\frac{\bar\cN}{\Lambda_1}\frac s{y_\ast^2}
\leq\frac{\bar\varpi\,s}{8\pi y_\ast}\leq\frac1{64}$, and
Lemma~\ref{lem:estimates}(iv) gives
$\norm{\delta_{i_0}}^2\leq\frac{16\pi^2}{y_\ast^2}\frac{q_{\mathrm{out}}s}\lambda
\leq\frac{2\pi\bar\varpi}{y_\ast}s$.
\end{proof}

\begin{proposition}[No arrow out of $i_0$]\label{prop:sink}
Let $\mathsf E$ have a surjective differential at $i_0$ and $r\geq1$, and
let $\rho_a=0$ for every arrow $a$ out of $i_0$; this holds if no arrow leaves
$i_0$, and if $W_j=0$ for all $j\neq i_0$. Then, for every metric, the complex is the orthogonal
direct sum of its cohomology representation and of the complex
$(\ker D)^\perp\to U$ at $i_0$, the flow preserves this decomposition,
$\delta_i=0$ for all $i$, the phases on $(\ker D)^\perp$ and $U$ are
$\alpha_\ast$ and $\alpha_\ast-\pi$, and
\[
 F(t)=e^{\pi t/2}F(0),\qquad
 \Omega_V=(z_\ast+\mu_{i_0})\oplus z_\ast e^{-F^\dagger F},\qquad
 \Omega_U=-z_\ast e^{-L}
\]
for all $t$. In particular
$\norm{\Omega_V-(z_\ast+\mu_{i_0})\oplus0}_1+\norm{\Omega_U}_1
\leq2|z_\ast|r\exp(-\lambda_0e^{\pi t})$, and the exact part decays. This
bound is at most $2|z_\ast|r\,e^{-(t-t_1)}$ for
$t\geq t_1:=\frac1\pi\log^+\frac1{\lambda_0}$.
\end{proposition}

\begin{proof}
Here $C_{\mathrm{out}}=0$, so $J=0$, $E_R=0$, $C_T=0$ and $Y_j=0$ in
\eqref{eq:blocks} and \eqref{eq:adaptedblocks}. The blocks of $\Omega_V$ are
$z_\ast+\mu_{i_0}$ and $z_\ast e^{-F^\dagger F}$, so $\Phi_{HA}=0$,
$\Phi_{AA}=\alpha_\ast$, $\Phi_U=\alpha_\ast-\pi$, and \eqref{eq:adapted} gives
$\dot F=\frac\pi2F$.
\end{proof}

\begin{remark}\label{rem:decay}
(1) The time $t_1$ contains the factor $e^{\Lambda_1}$, and
$\Lambda_1\geq8\bar\cN/y_\ast$.

(2) When an arrow leaves $i_0$, the phases $\Phi_{AA}$ and $\Phi_U$ on the
exact part need not tend to $\alpha_\ast$ and $\alpha_\ast-\pi$, in contrast
with Proposition~\ref{prop:sink}. In the proof of Theorem~\ref{thm:decay}(i), the bound
$\inner v{\Phi_Vv}-\inner u{\Phi_Uu}\geq\pi/2$ is used at a possible
contact with the line $\lambda=\frac{\bar\varpi}4(t-t_1)-\mathfrak c_3$.
It gives the stated lower bound for $\lambda(t)$.
\end{remark}

\section{Stability, the Kempf--Ness function and convergence}\label{sec:convergence}

\subsection{The stability parameter}
Let $\mathcal H$ be a nonzero representation of $Q$ with dimension vector
$(n_i)$ and $Z(\mathcal H)=|Z|e^{\ii\varphi}$, $\varphi\in(0,\pi)$,
$\sigma=\sin\varphi$; in the applications it is the cohomology of a complex
with a surjective differential, or the representation itself. For $i\in Q_0$
put
\[
 \theta_i=\frac{\Ima(e^{-\ii\varphi}z_i)}\sigma=\frac{|z_i|\sin(\alpha_i-\varphi)}\sigma .
\]
Then $\sum_in_i\theta_i=0$, and for a real number $x$ one has
$\arg(z_i+x)=\varphi$ if and only if $x=\theta_i$; moreover
$|z_i+\theta_i|=y_i/\sigma$. A representation $\mathsf R$ has
$\arg Z(\mathsf R)\leq\varphi$ if and only if
$\sum_i\theta_i\dim\mathsf R_i\leq0$. A representation $\mathsf R$ with
$\sum_i\theta_i\dim\mathsf R_i=0$, such as $\mathcal H$, is \emph{semistable}
if $\sum_i\theta_i\dim\mathsf R'_i\leq0$ for each of its subrepresentations
$\mathsf R'$, \emph{stable} if the inequality is strict for the proper
nonzero ones, and \emph{polystable} if it is a direct sum of stable
representations $\mathsf R''$ with $\sum_i\theta_i\dim\mathsf R''_i=0$. This
is stability in the sense of King \cite[Definition~1.1]{kin94} for the real
parameter $-\theta$ (the sign comes from the convention, outgoing minus
incoming, for the moment map), and it is the semistability of objects of
phase $\varphi/\pi$ for the stability condition of
Section~\ref{sec:setting}. Since $\arg(z_i+x)=\varphi$ only for
$x=\theta_i$, a metric on $\mathcal H$ is harmonic exactly when $\mu_i=\theta_i$ for all $i$; we call such a metric
a \emph{zero of the shifted moment map} $\mu-\theta$, where $\mu=(\mu_i)$ is
the moment map of Section~\ref{sec:setting}.

\subsection{The Kempf--Ness function}
Let $\mathcal X$ be the space of metrics
$\mathfrak h=(\mathfrak h_i)_{i\in Q_0}$ on the spaces $\mathcal H_i$. A
tangent vector at $\mathfrak h$ is a tuple $\xi=(\xi_i)$ of
$\mathfrak h$-Hermitian endomorphisms, $\delta\mathfrak h=\mathfrak h\xi$, and
$\norm\xi^2=\sum_i\tr\xi_i^2$ defines an invariant Riemannian metric for which
$\mathcal X$ is complete, simply connected and nonpositively curved, with
geodesics $u\mapsto\mathfrak he^{u\xi}$ and distance $d$. The
\emph{Kempf--Ness function} of $\mathcal H$ is
\begin{equation}\label{eq:KN}
 \cK_{\mathcal H}(\mathfrak h)=\sum_a\tr\bigl(\mathfrak h_{\mathsf s(a)}^{-1}\beta_a^*\,\mathfrak h_{\mathsf t(a)}\beta_a\bigr)
 +\sum_i\theta_i\log\det\mathfrak h_i ,
\end{equation}
where $\beta_a$ are the matrices of the arrows of $\mathcal H$ in fixed bases,
$\beta_a^*$ their conjugate transposes, and $\mathfrak h_i$ stands for the
Gram matrix. The first sum is the sum of the squared norms of the arrows for
the metric $\mathfrak h$. Since $\sum_in_i\theta_i=0$, $\cK_{\mathcal H}$ is
invariant under a common rescaling of the $\mathfrak h_i$. In terms of the
function \eqref{eq:primitive} for the representation $\mathcal H$,
$\cK_{\mathcal H}=\Ima(e^{-\ii\varphi}\cS_{\CC})/\sigma$.

\begin{lemma}\label{lem:variation}
The differential of the Kempf--Ness function is
$d\cK_{\mathcal H}(\xi)=-\sum_i\tr\bigl(\xi_i(\mu_i-\theta_i)\bigr)$, and along
a geodesic
\[
 \frac{d^2}{du^2}\cK_{\mathcal H}(\mathfrak he^{u\xi})
 =\sum_a\norm{\xi_{\mathsf t(a)}\beta_a(u)-\beta_a(u)\xi_{\mathsf s(a)}}^2\geq0,\qquad
 \beta_a(u)=e^{u\xi_{\mathsf t(a)}/2}\beta_ae^{-u\xi_{\mathsf s(a)}/2},
\]
in an $\mathfrak h$-orthonormal frame. So $\cK_{\mathcal H}$ is geodesically
convex, its critical points are the zeros of the shifted moment map, and its
gradient has the norm $\norm{\mu-\theta}:=(\sum_i\norm{\mu_i-\theta_i}^2)^{1/2}$.
\end{lemma}

\begin{proof}
In an $\mathfrak h$-orthonormal frame the first sum of \eqref{eq:KN} at
$\mathfrak he^{u\xi}$ is
$\sum_a\tr(e^{-u\xi_{\mathsf s(a)}}\beta_a^\dagger e^{u\xi_{\mathsf t(a)}}\beta_a)$,
and $\log\det$ is affine along geodesics with derivative $\tr\xi_i$.
Differentiate once and twice.
\end{proof}

For $\mathfrak h\in\mathcal X$ let $\mathfrak a(\mathfrak h)$ be the space of
tuples $\xi$ of $\mathfrak h$-Hermitian endomorphisms with
$\xi_{\mathsf t(a)}\beta_a=\beta_a\xi_{\mathsf s(a)}$ for all arrows $a$, the
Hermitian infinitesimal automorphisms of $\mathcal H$, and let
$\mathcal Z\subset\mathcal X$ be the set of zeros of the shifted moment map. The next
lemma collects the convexity properties used; they are classical for
Kempf--Ness functions (see \cite[Theorem~4.3]{grs21} for actions on closed
K\"ahler manifolds), and we include the short proofs for the present linear
action.

\begin{lemma}\label{lem:convexity}
Suppose $\mathcal Z\neq\emptyset$.
\begin{enumerate}
\item[(a)] For $p\in\mathcal Z$, $\mathcal Z=\{pe^\xi:\xi\in\mathfrak a(p)\}$;
 it is a closed totally geodesic submanifold.
\item[(b)] The number
 \[
 \nu_*=\min\Bigl\{\sum_a\norm{\xi_{\mathsf t(a)}\beta_a-\beta_a\xi_{\mathsf s(a)}}^2:\
 \xi\perp\mathfrak a(p),\ \norm\xi=1\Bigr\}>0,
 \]
 with the $\beta_a$ in a $p$-orthonormal frame, does not depend on
 $p\in\mathcal Z$ (if $\mathfrak a(p)$ is the whole tangent space, which
 happens when all arrows vanish, put $\nu_*=\infty$, with the convention
 $\infty\cdot0=0$ below; then $\mathcal Z=\mathcal X$ and
 $\cK_{\mathcal H}$ is constant); and for every unit
 $\xi\perp\mathfrak a(p)$ and
 $0\leq u\leq\varrho_*:=\log2/\sqrt2$,
 $\frac{d^2}{du^2}\cK_{\mathcal H}(pe^{u\xi})\geq\nu_*/2$.
\item[(c)] For $\mathfrak h\in\mathcal X$ put
 $\mathfrak d=d(\mathfrak h,\mathcal Z)$ and
 $\mathcal E=\cK_{\mathcal H}(\mathfrak h)-\min\cK_{\mathcal H}$. Then
 $\mathcal E\leq\norm{\mu-\theta}\,\mathfrak d$,
 $\mathcal E\geq\frac{\nu_*}4\min(\mathfrak d^2,\varrho_*\mathfrak d)$,
 and if $\mathfrak d\leq\varrho_*$, then
 $\norm{\mu-\theta}\geq\frac{\nu_*}2\mathfrak d$ and
 $\mathcal E\leq\frac2{\nu_*}\norm{\mu-\theta}^2$.
\end{enumerate}
\end{lemma}

\begin{proof}
(a) If $\xi\in\mathfrak a(p)$, the second derivative of $\cK_{\mathcal H}$
along $u\mapsto pe^{u\xi}$ vanishes identically by
Lemma~\ref{lem:variation}, and the first derivative vanishes at $u=0$; so
$pe^\xi\in\mathcal Z$. Conversely, if $pe^\xi\in\mathcal Z$, the convex
function $u\mapsto\cK_{\mathcal H}(pe^{u\xi})$ takes its minimum at $u=0$ and
$u=1$, hence is constant on $[0,1]$, and its second derivative at $0$ gives
$\xi\in\mathfrak a(p)$. So $\mathcal Z$ is the image of a linear subspace
under the exponential map at each of its points.

(b) The quadratic form is positive on the orthogonal complement of its kernel
$\mathfrak a(p)$. By Lemma~\ref{lem:variation},
$\xi_{\mathsf t(a)}\beta_a(u)-\beta_a(u)\xi_{\mathsf s(a)}
=e^{u\xi_{\mathsf t(a)}/2}(\xi_{\mathsf t(a)}\beta_a-\beta_a\xi_{\mathsf s(a)})e^{-u\xi_{\mathsf s(a)}/2}$
has norm at least
$e^{-u(\norm{\xi_{\mathsf t(a)}}_{\mathrm{op}}+\norm{\xi_{\mathsf s(a)}}_{\mathrm{op}})/2}$
times that of $\xi_{\mathsf t(a)}\beta_a-\beta_a\xi_{\mathsf s(a)}$, and
$\norm{\xi_{\mathsf s(a)}}_{\mathrm{op}}+\norm{\xi_{\mathsf t(a)}}_{\mathrm{op}}\leq\sqrt2$
for a unit $\xi$, because $\mathsf s(a)\neq\mathsf t(a)$. So the second
derivative is at least $e^{-\sqrt2u}\nu_*$. If $p'=pe^\zeta$ with
$\zeta\in\mathfrak a(p)$, then $e^{-\zeta/2}$ is a tuple of automorphisms of
$\mathcal H$ that maps $p$-orthonormal frames to $p'$-orthonormal frames and
preserves the matrices $\beta_a$; conjugation by it carries $\mathfrak a(p)$
and its orthogonal complement isometrically to those of $p'$.

(c) Let $p\in\mathcal Z$ be a nearest point and
$\mathfrak h=pe^{\mathfrak d\xi}$ with $\norm\xi=1$; then
$\xi\perp\mathfrak a(p)$. The function
$f(u)=\cK_{\mathcal H}(pe^{u\xi})-\min\cK_{\mathcal H}$ is convex with
$f(0)=f'(0)=0$. First, $0=f(0)\geq f(\mathfrak d)-\mathfrak df'(\mathfrak d)$
and $f'(\mathfrak d)\leq\norm{\mu-\theta}$. Secondly,
$f''\geq\nu_*/2$ on $[0,\varrho_*]$ by (b), so $f(u)\geq\frac{\nu_*}4u^2$
there, and $f(u)\geq f(\varrho_*)u/\varrho_*$ for $u\geq\varrho_*$ by
convexity. Thirdly, if $\mathfrak d\leq\varrho_*$, then
$\norm{\mu-\theta}\geq f'(\mathfrak d)=\int_0^{\mathfrak d}f''\geq\frac{\nu_*}2\mathfrak d$;
combine with the first inequality.
\end{proof}

\begin{lemma}[Kempf--Ness, King]\label{lem:king}
Let $\cK_*=\inf\cK_{\mathcal H}\in[-\infty,\infty)$.
\begin{enumerate}
\item[(a)] $\cK_*>-\infty$ if and only if $\mathcal H$ is semistable.
\item[(b)] A zero of the shifted moment map exists if and only if $\mathcal H$ is
 polystable.
\end{enumerate}
\end{lemma}

For a parameter that is a positive multiple of an integral vector these are
the theorems of Kempf--Ness and King \cite{kn79,kin94}: (b) is
\cite[Proposition~6.5]{kin94}, and (a) is stated in
\cite[Theorem~5.1]{hkkp23} and follows from
\cite[Proposition~3.1, Lemmas~2.2 and~6.3]{kin94}. The parameter $\theta$ is
real, and for a quiver with more than two vertices it need not be a multiple
of an integral vector; (b) is stated for real parameters in
\cite[Theorem~5.1]{hkkp26}, with reference to \cite{kin94}, and the general
theory of \cite{grs21} has no integrality hypothesis but is written for
actions on closed K\"ahler manifolds. We give the reduction to integral
parameters.

\begin{proof}
Let $\Theta$ be the real vector space of all $\theta'\in\RR^{Q_0}$ with
$\sum_in_i\theta_i'=0$, and write $\cK^{\theta'}$ for the function
\eqref{eq:KN} with $\theta'$ in the place of $\theta$. The first sum of
\eqref{eq:KN} does not depend on $\theta'$, so
$\cK^{\sum c_l\theta^{(l)}}=\sum c_l\cK^{\theta^{(l)}}$ whenever $\sum c_l=1$. For
a positive number $c$, $\cK^{c\theta'}$ at the arrows $\beta$ is $c$ times
$\cK^{\theta'}$ at the arrows $\beta/\sqrt c$, which have the same
subrepresentations; so the cited theorems hold for all rational
$\theta'\in\Theta$. Let $\mathfrak S$ be the finite set of dimension vectors of
subrepresentations of $\mathcal H$.

(a) If a subrepresentation $\mathcal H'$ with dimension vector $(n_i')$ has
$\sum_in_i'\theta_i>0$, let $\mathfrak h_i=1\oplus e^\tau$ on
$\mathcal H_i'\oplus\mathcal H_i'^\perp$ for reference metrics. The first sum
of \eqref{eq:KN} stays bounded as $\tau\to\infty$, because the arrows map
$\mathcal H'$ to $\mathcal H'$, and
$\sum_i\theta_i\log\det\mathfrak h_i=\tau\sum_i\theta_i(n_i-n_i')=-\tau\sum_in_i'\theta_i\to-\infty$.
Conversely, the set of $\theta'\in\Theta$ for which $\mathcal H$ is
semistable is the rational polyhedral cone
$\{\theta':\sum_in_i'\theta_i'\leq0\text{ for all }(n_i')\in\mathfrak S\}$. A point
of a rational polyhedral cone is a convex combination of rational points of
the cone: the cone is generated by finitely many rational vectors $v_j$, and
if $\theta=\sum_j\lambda_jv_j$ with $\lambda_j\geq0$, choose rational numbers
$q_j>0$ with $\sum_j\lambda_j/q_j<1$; then $\theta$ is a convex combination of
the points $q_jv_j$ and $0$. So $\theta=\sum_lc_l\theta^{(l)}$ with $c_l\geq0$, $\sum c_l=1$ and
rational $\theta^{(l)}$ for which $\mathcal H$ is semistable, and
$\cK^\theta=\sum_lc_l\cK^{\theta^{(l)}}$ is bounded below.

(b) Let $\mathfrak h$ be a zero of the shifted moment map, $\mathcal H'$ a
subrepresentation and $\Pi_i$ the $\mathfrak h$-orthogonal projection onto
$\mathcal H_i'$. Since
$\beta_a\Pi_{\mathsf s(a)}=\Pi_{\mathsf t(a)}\beta_a\Pi_{\mathsf s(a)}$,
\[
 \sum_in_i'\theta_i=\sum_i\tr(\Pi_i\mu_i)
 =\sum_a\bigl(\norm{\beta_a\Pi_{\mathsf s(a)}}^2-\norm{\Pi_{\mathsf t(a)}\beta_a}^2\bigr)
 =-\sum_a\norm{\Pi_{\mathsf t(a)}\beta_a(1-\Pi_{\mathsf s(a)})}^2\leq0 .
\]
So $\mathcal H$ is semistable, and if equality holds the orthogonal
complement of $\mathcal H'$ is a subrepresentation, and $\mathfrak h$
restricts to a zero of the shifted moment map on both summands. By induction on the
dimension, $\mathcal H$ is polystable. Conversely, an orthogonal direct sum
of zeros of the shifted moment maps of the summands is a zero of the shifted moment map; so
let $\mathcal H$ be stable. The set of $\theta'\in\Theta$ for which
$\mathcal H$ is stable is open and defined by finitely many strict rational
inequalities, so $\theta=\sum_lc_l\theta^{(l)}$ with $c_l>0$, $\sum c_l=1$ and
rational $\theta^{(l)}$ for which $\mathcal H$ is stable. For each $l$ the
function $\cK^{\theta^{(l)}}$ attains its minimum
\cite[Proposition~6.5]{kin94}, and by Lemma~\ref{lem:convexity}(a), which
holds for every parameter in $\Theta$, its set of minimizers is
$\{pe^\xi:\xi\in\mathfrak a(p)\}$; a stable representation has only scalar
endomorphisms, so the minimizers form one orbit of the common rescalings. A
geodesically convex function on the quotient of $\mathcal X$ by the common
rescalings, which is again a complete simply connected manifold of
nonpositive curvature, with a single minimizer is proper: it exceeds its
minimum by a positive constant on the unit sphere around the minimizer, and
then by that constant times $d$ at distance $d\geq1$. Hence
$\cK^\theta=\sum_lc_l\cK^{\theta^{(l)}}$ is proper on this quotient and attains
its minimum, at a critical point.
\end{proof}

\subsection{Convergence}
For $i\in Q_0$ let $G_i(u)=\varphi-\arg(z_i+u)$ for real $u$. These functions
are increasing, $G_i(\theta_i)=0$, and
$0<G_i'(u)=y_i/|z_i+u|^2\leq1/y_i$. For a number $\bar\cN\geq0$ put
\begin{equation}\label{eq:Gamma}
 \bar Q=\max\bigl\{\bar\cN,\ \max_i|\theta_i|\bigr\},\qquad
 \Gamma=\max_i\frac{(|z_i|+\bar Q)^2}{y_i}.
\end{equation}
The next proposition is the common part of the convergence theorems for
representations and for complexes; for representations the perturbation
$\delta$ in it is zero.

\begin{proposition}\label{prop:perturbed}
Let $\mathfrak h(t)$, $t\geq\tau_0$, be a differentiable curve in $\mathcal X$
whose velocity $\xi=\mathfrak h^{-1}\dot{\mathfrak h}$ has the form
\[
 \xi_i=G_i(\mu_i)-\delta_i\qquad(i\in Q_0),
\]
with $\mu=\mu(\mathfrak h(t))$ and $\mathfrak h$-Hermitian endomorphisms
$\delta_i$. Put $\norm\delta^2=\sum_i\norm{\delta_i}^2$, and suppose that for
$t\geq\tau_0$ the sum of the squared norms of the arrows for the metric
$\mathfrak h(t)$ is at most $\bar\cN$, and that
$\norm{\delta(t)}\leq c_\delta\,e^{-\kappa(t-\tau_0)}$ with constants
$c_\delta\geq0$, $\kappa>0$. Let $\Gamma$ be as in \eqref{eq:Gamma}. Then:
\begin{enumerate}
\item[(i)] $\frac d{dt}\cK_{\mathcal H}(\mathfrak h(t))
 \leq-\frac1\Gamma\norm{\mu-\theta}^2+\norm{\mu-\theta}\,\norm\delta$ for
 $t\geq\tau_0$.
\item[(ii)] If $\mathcal H$ is semistable, then $\mu_i\to\theta_i$ for all
 $i$ and $\cK_{\mathcal H}(\mathfrak h(t))\to\cK_*$ as $t\to\infty$.
\item[(iii)] If $\mathcal H$ is polystable, there is a zero
 $\mathfrak h_\infty$ of the shifted moment map with
 $d(\mathfrak h(t),\mathfrak h_\infty)\leq\mathfrak C\,e^{-\varkappa t}$ for
 every $\varkappa<\min\{\nu_*/(8\Gamma),\kappa\}$, with a constant
 $\mathfrak C$ depending on $\varkappa$ and the curve; and $\mu-\theta$ tends
 to zero exponentially.
\end{enumerate}
\end{proposition}

\begin{proof}
(i) By Lemma~\ref{lem:variation},
\[
 \frac d{dt}\cK_{\mathcal H}(\mathfrak h(t))=-\sum_i\tr\bigl((\mu_i-\theta_i)G_i(\mu_i)\bigr)
 +\sum_i\tr\bigl((\mu_i-\theta_i)\delta_i\bigr).
\]
The eigenvalues of $\mu_i$ lie in $[-\bar\cN,\bar\cN]$; for real $u$ with
$|u|\leq\bar Q$, $(u-\theta_i)G_i(u)=G_i'(\zeta)(u-\theta_i)^2$ with a point
$\zeta$ between $u$ and $\theta_i$, and $G_i'(\zeta)\geq1/\Gamma$.

(ii) Put $\mathcal E(t)=\cK_{\mathcal H}(\mathfrak h(t))-\cK_*\geq0$
(Lemma~\ref{lem:king}(a)). By (i),
\[
 \frac d{dt}\cK_{\mathcal H}(\mathfrak h(t))\leq-\frac1{2\Gamma}\norm{\mu-\theta}^2+\frac\Gamma2\norm\delta^2 ,
\]
so $\norm{\mu-\theta}^2$ is integrable on $[\tau_0,\infty)$. In
$\mathfrak h$-orthonormal frames the arrows evolve by
$\dot\beta_a=\frac12(\xi_{\mathsf t(a)}\beta_a-\beta_a\xi_{\mathsf s(a)})$, and
$\beta_a$ and $\xi$ are bounded; so the derivatives of the $\mu_i$ are
bounded, $\norm{\mu-\theta}^2$ is uniformly continuous, and
$\mu_i\to\theta_i$. The function
$\mathcal E(t)+\frac\Gamma2\int_t^\infty\norm\delta^2$ is nonincreasing, so
$\mathcal E$ has a limit $\mathcal E_\infty\geq0$. Suppose
$\mathcal E_\infty>0$ and choose $\mathfrak h'$ with
$\cK_{\mathcal H}(\mathfrak h')-\cK_*<\mathcal E_\infty/2$. By convexity along
the geodesic from $\mathfrak h(t)$ to $\mathfrak h'$,
$\cK_{\mathcal H}(\mathfrak h(t))-\cK_{\mathcal H}(\mathfrak h')\leq\norm{\mu-\theta}\,d(\mathfrak h(t),\mathfrak h')$,
and the left side is at least $\mathcal E_\infty/4$ for large $t$. The speed
of $\mathfrak h(t)$ is at most $\norm{\mu-\theta}/\min_iy_i+\norm\delta$,
which is square integrable, so
$d(\mathfrak h(t),\mathfrak h')\leq c+c'\sqrt t$; then
$\norm{\mu-\theta}^2\geq\mathcal E_\infty^2/(16(c+c'\sqrt t)^2)$ for large
$t$, which is not integrable. So $\mathcal E\to0$.

(iii) By Lemma~\ref{lem:king}(b), $\mathcal Z\neq\emptyset$. By
Lemma~\ref{lem:convexity}(c) and $\mathcal E\to0$,
$d(\mathfrak h(t),\mathcal Z)\leq\varrho_*$ for large $t$, and then
$\norm{\mu-\theta}^2\geq\frac{\nu_*}2\mathcal E$, so
$\frac d{dt}\mathcal E\leq-\frac{\nu_*}{4\Gamma}\mathcal E+\frac\Gamma2\norm\delta^2$
and $\mathcal E(t)=O(e^{-2\varkappa t})$ for every
$2\varkappa<\min\{\nu_*/(4\Gamma),2\kappa\}$. Integrating the second
inequality of the proof of (ii) over $[t,t+1]$ bounds
$\int_t^{t+1}\norm{\mu-\theta}^2$ by a multiple of
$\mathcal E(t)+\int_t^{t+1}\norm\delta^2$; by the bound on the speed and the
Cauchy--Schwarz inequality the length of $\mathfrak h$ over $[t,t+1]$ is
$O(e^{-\varkappa t})$, so $\mathfrak h(t)$ converges at this rate, and the
limit is a zero of the shifted moment map by continuity. Finally $\mu$ is a
locally Lipschitz function on $\mathcal X$ with $\mu(\mathfrak h_\infty)=\theta$.
\end{proof}

\begin{theorem}[Representations]\label{thm:representations}
Let $Q$ be a finite quiver without oriented cycles, with charges $z_i$ in the
upper half-plane, and let $\mathsf R$ be a nonzero representation of $Q$
with $Z(\mathsf R)=|Z|e^{\ii\varphi}$ and
$\theta_i=\Ima(e^{-\ii\varphi}z_i)/\sin\varphi$. Let $h(t)$ be the solution
of the phase flow $h_i^{-1}\dot h_i=-\arg(z_i+\mu_i)$ from an arbitrary
initial metric, and $\mathfrak h(t)=e^{\varphi t}h(t)$. Then:
\begin{enumerate}
\item[(a)] the Kempf--Ness function $\cK_{\mathsf R}(\mathfrak h(t))=\cK_{\mathsf R}(h(t))$ is
 nonincreasing;
\item[(b)] if $\mathsf R$ is semistable, then $\mu_i(t)\to\theta_i$ for all
 $i$, $\cK_{\mathsf R}(h(t))\to\inf\cK_{\mathsf R}$, and the mass
 $\sum_i\tr|z_i+\mu_i(t)|$ tends to $|Z|$;
\item[(c)] if $\mathsf R$ is polystable, $\mathfrak h(t)$ converges to a
 harmonic metric $\mathfrak h_\infty$, that is, a metric with
 $\mu_i(\mathfrak h_\infty)=\theta_i$ for all $i$, and
 $d(\mathfrak h(t),\mathfrak h_\infty)\leq\mathfrak C\,e^{-\varkappa t}$ for
 every $\varkappa<\nu_*/(8\Gamma)$, where $\nu_*$ is the number of
 Lemma~\ref{lem:convexity}(b) and $\Gamma$ is given by \eqref{eq:Gamma} with
 $\bar\cN=\bar\cN_{\mathsf R}$, the bound of
 Proposition~\ref{prop:reparrows}; and $\mu-\theta$ and the mass minus $|Z|$
 tend to zero exponentially.
\end{enumerate}
\end{theorem}

\begin{proof}
The block of $\Omega$ at $i$ is $z_i+\mu_i$, so the velocity of
$\mathfrak h(t)$ is $\xi_i=\varphi-\arg(z_i+\mu_i)=G_i(\mu_i)$. By
Lemma~\ref{lem:variation},
$\frac d{dt}\cK_{\mathsf R}(\mathfrak h(t))=-\sum_i\tr\bigl((\mu_i-\theta_i)G_i(\mu_i)\bigr)\leq0$,
because $(u-\theta_i)G_i(u)\geq0$ for all real $u$; this is (a). By
Proposition~\ref{prop:reparrows} the sum of the squared norms of the arrows
is at most $\bar\cN_{\mathsf R}$ for $t\geq t_{\mathsf R}$, and it does not
change under a common rescaling of the metrics. So
Proposition~\ref{prop:perturbed} applies with $\tau_0=t_{\mathsf R}$,
$\delta=0$ and every $\kappa>0$, and gives (b) and (c); for the mass,
$\sum_i\tr|z_i+\mu_i|\to\sum_in_i|z_i+\theta_i|=\sum_in_iy_i/\sigma=|Z|$,
exponentially in case (c) because $\tr|z_i+\mu_i|$ is a Lipschitz function of
$\mu_i$.
\end{proof}

\begin{theorem}[Complexes with a surjective differential]\label{thm:convergence}
Let $\mathsf E$ be a complex with a surjective differential at $i_0$ and
$r\geq1$, with $W_j\neq0$ for some $j\neq i_0$ and with at least one arrow
leaving $i_0$; fix a height function and let the constants be as in
\eqref{eq:constants}, and $\Gamma$ as in \eqref{eq:Gamma}. Let
\[
 \mathfrak h(t)=\bigl(e^{\varphi t}h_i(t)|_{\mathcal H_i}\bigr)_{i\in Q_0}
\]
be the metrics induced on the cohomology representation $\mathcal H$,
multiplied by $e^{\varphi t}$; at $i_0$ the induced metric is the restriction
of the metric of $V$ to $\ker D$. Put
$c_\delta^2=\frac{2\pi\bar\varpi\,s_1}{y_\ast}+s_1^2\sum_{j\neq i_0}y_j^{-2}$.
For all charges in the upper half-plane:
\begin{enumerate}
\item[(i)] For $t\geq t_2$,
 \[
 \frac d{dt}\cK_{\mathcal H}(\mathfrak h(t))\leq-\frac1\Gamma\norm{\mu-\theta}^2+\norm{\mu-\theta}\,\norm\delta,
 \qquad\norm{\delta(t)}\leq c_\delta\,e^{-\frac{\bar\varpi}4(t-t_1)}.
 \]
\item[(ii)] If $\mathcal H$ is semistable, then as $t\to\infty$:
 $\mu_i\to\theta_i$ for all $i$, $\cK_{\mathcal H}(\mathfrak h(t))\to\cK_*$
 and $\mass(h(t))\to|Z|$.
\item[(iii)] If $\mathcal H$ is polystable, $\mathfrak h(t)$ converges to a
 harmonic metric $\mathfrak h_\infty$ of $\mathcal H$, a zero of the shifted
 moment map, with
 $d(\mathfrak h(t),\mathfrak h_\infty)\leq\mathfrak C\,e^{-\varkappa t}$ for
 every $\varkappa<\min\{\nu_*/(8\Gamma),\ \bar\varpi/4\}$; and $\mu-\theta$ and
 $\mass-|Z|$ tend to zero exponentially.
\end{enumerate}
\end{theorem}

\begin{proof}
By \eqref{eq:flow} the velocity of $\mathfrak h(t)$ at the vertex $i$ is
$\varphi-\Phi_i$, with $\Phi_{HH}$ in the place of $\Phi_{i_0}$, that is,
$\xi_i=G_i(\mu_i)-\delta_i$ with the $\delta_i$ of
Section~\ref{sec:setting}. For $t\geq t_0$ the sum of the squared norms of
the arrows of $\mathcal H$ is at most $\cN\leq\bar\cN$ by
\eqref{eq:energybounds} and Theorem~\ref{thm:energy}, and for $t\geq t_2$
Theorem~\ref{thm:decay}(i),(iii) gives
$\norm\delta^2\leq\frac{2\pi\bar\varpi}{y_\ast}s+s^2\sum_jy_j^{-2}\leq c_\delta^2e^{-\frac{\bar\varpi}2(t-t_1)}$.
So Proposition~\ref{prop:perturbed} applies with $\tau_0=t_2$ and
$\kappa=\bar\varpi/4$. For the mass,
$\sum_i\norm{z_i+\mu_i}_1\to\sum_in_i|z_i+\theta_i|=|Z|$, exponentially in
case (iii) because $\norm{z_i+\mu_i}_1$ is a Lipschitz function of $\mu_i$,
and the difference to $\mass$ is bounded by Theorem~\ref{thm:decay}(ii).
\end{proof}

\begin{remark}\label{rem:vanishing}
Let $\mathsf E$ have a surjective differential at $i_0$ and $r\geq1$, and
let $\rho_a=0$ for all arrows $a$ out of $i_0$; this includes the case in
which all $W_j$ vanish. By Proposition~\ref{prop:sink} the complex is, for
every metric, the orthogonal direct sum of its cohomology representation
and an exact complex at $i_0$, and the flow preserves this decomposition. So
the metrics induced on the cohomology follow the phase flow of the
representation $\mathcal H$, Theorem~\ref{thm:representations} applies to
them, and $\mass(h(t))\to|Z|$ if $\mathcal H$ is semistable, exponentially if
it is polystable.
\end{remark}

For example, Theorems~\ref{thm:decay} and~\ref{thm:convergence},
Proposition~\ref{prop:sink} and Remark~\ref{rem:vanishing} apply to the
quiver $A_l$ with the linear orientation, with the differential at any
vertex.

\subsection{Complexes with an injective differential}
Let $\mathsf E'$ be a complex with an injective differential at $i_0$
(Definition~\ref{def:type}), with $U'\xrightarrow{D'}V'$ at $i_0$. The
arrows out of $i_0$ vanish on $D'U'$, and the cohomology is a representation
in degree zero with $V'/D'U'$ at $i_0$; the \emph{exact part} is the
subcomplex $U'\to D'U'$, of rank $\dim U'$. The metric induced on $V'/D'U'$ is
the quotient metric, that is, the metric of the orthogonal complement of
$D'U'$ in $V'$.

\begin{proposition}[Injective differential]\label{prop:injective}
Let $\mathsf E'$ be a complex with an injective differential at $i_0$ and
$U'\neq0$, with total phase $\varphi$, and let $\lambda'$ be the least
eigenvalue of $D'^\dagger D'$. Its dual $\mathsf E'^\vee$ over the opposite quiver, with
the charges $-\overline{z_i}$, is a complex with a surjective differential at
$i_0$ and total phase $\pi-\varphi$. Along the phase flow of $\mathsf E'$:
\begin{enumerate}
\item[(i)] $\lambda'\to\infty$, and the blocks of $\Omega$ approach those of
 the cohomology representation with its induced metrics, with the bound of
 Theorem~\ref{thm:decay}(ii), or of Proposition~\ref{prop:sink}, for
 $\mathsf E'^\vee$;
\item[(ii)] the statements of Theorem~\ref{thm:convergence}, or of
 Remark~\ref{rem:vanishing}, hold for the metrics induced on the cohomology
 of $\mathsf E'$, multiplied by $e^{\varphi t}$, with the constants of
 $\mathsf E'^\vee$.
\end{enumerate}
\end{proposition}

\begin{proof}
The opposite quiver has no oriented cycles, and the transposed differential
$D'^{\mathsf T}:V'^*\to U'^*$ is surjective in degrees $0,1$. By
Proposition~\ref{prop:usedc}(b), the trajectories correspond by
$h\mapsto e^{-\pi t}(h^{-1})^{\mathsf T}$, the blocks of $\Omega$ correspond
up to sign and complex conjugation in dual orthonormal frames, with the same
trace norms, and the mass is preserved. The cohomology representation of
$\mathsf E'^\vee$ is the dual of that of $\mathsf E'$, with the dual of the
induced metrics; subrepresentations of one correspond to quotients of the
other, and the argument $\psi$ of a central charge corresponds to
$\pi-\psi$, so that one is semistable, stable or polystable if and only if
the other is. The normalized induced metrics
$e^{(\pi-\varphi)t}e^{-\pi t}(h^{-1})^{\mathsf T}$ of $\mathsf E'^\vee$
converge exactly when $e^{\varphi t}h$ does.
\end{proof}

\section{Semistable and unstable cohomology}\label{sec:asymptotics}

This section describes the long-time behaviour when $\mathcal H$ is not
polystable. For semistable $\mathcal H$ the curve of normalized metrics
stays at bounded distance from a trajectory of a gradient flow of the
Kempf--Ness function, whose asymptotics are known
(Theorem~\ref{thm:comparison}). With no hypothesis on $\mathcal H$, the mass
converges to the mass of the Harder--Narasimhan filtration, the filtration
of the spaces by the exponential rate of decay of the metric is that
filtration, and its factors, which are semistable, follow the semistable
theory (Theorem~\ref{thm:hn}).

Throughout the section $Q$ is a finite quiver without oriented cycles,
$\mathcal H$ is a nonzero representation of $Q$, with
$Z(\mathcal H)=|Z|e^{\ii\varphi}$, and $\mathfrak h(t)$,
$t\geq\tau_0$, is a curve of metrics on $\mathcal H$ as in
Proposition~\ref{prop:perturbed}: it is continuously differentiable, its
velocity is $\xi_i=G_i(\mu_i)-\delta_i$ with continuous
$\mathfrak h$-Hermitian $\delta_i$,
$\norm{\delta(t)}\leq c_\delta e^{-\kappa(t-\tau_0)}$ with $\kappa>0$, and the
sum of the squared norms of the arrows is at most $\bar\cN$. We call this
\emph{hypothesis} (P). No stability of $\mathcal H$ is assumed in (P). It
holds, with $\delta=0$, for the normalized phase flow
$\mathfrak h=e^{\varphi t}h$ on a representation of a quiver without oriented
cycles (Proposition~\ref{prop:reparrows}); and it holds for the normalized
metrics induced on the cohomology of a complex with a surjective or an
injective differential at one vertex, for every cohomology representation:
the proof of Theorem~\ref{thm:convergence}(i) and
Theorem~\ref{thm:decay} use no stability, and Remark~\ref{rem:vanishing}
covers the complexes excluded there. For an injective differential one
passes to the dual complex (Proposition~\ref{prop:injective}): the
normalized induced metrics of the two are dual to each other, so their
velocities are $\xi$ and $-\xi^{\mathsf T}$; the moment map changes sign
under duality; and
$\varphi-\arg(z_i+u)=-\bigl[(\pi-\varphi)-\arg(-\overline{z_i}-u)\bigr]$
for real $u$. So (P) for the dual gives (P).

We write $\Phi^{\mathcal H}_i=\arg(z_i+\mu_i)=\varphi-G_i(\mu_i)$ for the
phase of the representation $\mathcal H$ with the metric $\mathfrak h$, a
Hermitian matrix in an orthonormal frame; the function
$u\mapsto\arg(z_i+u)$ is decreasing, with derivative in $[-1/y_i,0)$. In
orthonormal frames of $\mathfrak h(t)$ whose coordinate maps $g_i$, with
$|g_iv|^2=\mathfrak h_i(v,v)$, satisfy $\dot g_ig_i^{-1}=\frac12\xi_i$, the
arrows evolve by
\begin{equation}\label{eq:betaflow}
 \dot\beta_a=\tfrac12\bigl(\beta_a(\Phi^{\mathcal H}_{\mathsf s(a)}+\delta_{\mathsf s(a)})
 -(\Phi^{\mathcal H}_{\mathsf t(a)}+\delta_{\mathsf t(a)})\beta_a\bigr),
\end{equation}
and the eigenvalues of all $\Phi^{\mathcal H}_i$ lie in a compact interval
$[\psi_-,\psi_+]\subset(0,\pi)$ determined by the charges and $\bar\cN$. Two
facts on functions of Hermitian matrices are used. If $g$ is a continuously
differentiable function on $\RR$ with $|g'|\leq c$, then
$\norm{g(A)-g(B)}\leq c\norm{A-B}$ for Hermitian $A,B$ and the
Hilbert--Schmidt norm; and if $g$ is nonincreasing, then
$\tr\bigl((g(A)-g(B))(A-B)\bigr)\leq0$. Both follow from the identities
$\inner{u_j}{(g(A)-g(B))v_l}=(g(a_j)-g(b_l))\inner{u_j}{v_l}$ and
$\inner{u_j}{(A-B)v_l}=(a_j-b_l)\inner{u_j}{v_l}$ for orthonormal
eigenvectors $u_j$ of $A$ and $v_l$ of $B$ with eigenvalues $a_j$, $b_l$.

\subsection{Semistable cohomology: comparison with the gradient flow}
Put
\[
 \mathsf m_i=\frac1{G_i'(\theta_i)}=\frac{y_i}{\sigma^2}\qquad(i\in Q_0),
\]
and let $d_{\mathsf m}$ be the distance of the Riemannian metric
$\norm\xi_{\mathsf m}^2=\sum_i\mathsf m_i\tr\xi_i^2$ on $\mathcal X$. With
this metric $\mathcal X$ is the product of the symmetric spaces of positive
definite Hermitian matrices, the $i$-th factor with its invariant metric
multiplied by $\mathsf m_i$; it is complete, simply connected and
nonpositively curved, its geodesics are again
$u\mapsto\mathfrak he^{u\xi}$, and
$d_{\mathsf m}(\mathfrak h,\mathfrak h')^2=\sum_i\mathsf m_i\norm{\log(\mathfrak h_i^{-1/2}\mathfrak h_i'\mathfrak h_i^{-1/2})}^2$.
By Lemma~\ref{lem:variation} the gradient of $\cK_{\mathcal H}$ for this
metric is the tangent vector $(-(\mu_i-\theta_i)/\mathsf m_i)_i$. The
\emph{weighted gradient flow} is
\begin{equation}\label{eq:gradflow}
 \mathfrak g_i^{-1}\dot{\mathfrak g}_i=\frac{\mu_i(\mathfrak g)-\theta_i}{\mathsf m_i}\qquad(i\in Q_0).
\end{equation}
Its solutions exist for all later times: the sum $q$ of the squared norms of
the arrows satisfies
$\dot q=-\sum_i\tr(\mu_i\xi_i)\leq\sum_in_i\theta_i^2/(4\mathsf m_i)$ along
\eqref{eq:gradflow}, so the velocity is bounded on bounded intervals of
time. The flow \eqref{eq:gradflow} is the gradient flow studied in
\cite{hkkp23}: it is \cite[(5.11)]{hkkp23}, with the parameters
$\mathsf m_i$ and $\theta_i$ there.

\begin{lemma}\label{lem:quadratic}
For all real $u$ and all $i$,
$\bigl|G_i(u)-(u-\theta_i)/\mathsf m_i\bigr|\leq\frac{3\sqrt3}{16\,y_i^2}(u-\theta_i)^2$.
Consequently
$\norm{G_i(M)-(M-\theta_i)/\mathsf m_i}\leq\frac{3\sqrt3}{16\,y_i^2}\norm{M-\theta_i}^2$
for every Hermitian matrix $M$.
\end{lemma}

\begin{proof}
$G_i''(u)=-2y_i(x_i+u)/((x_i+u)^2+y_i^2)^2$, and
$2y|w|/(w^2+y^2)^2$ has its maximum $3\sqrt3/(8y^2)$ at $w^2=y^2/3$. Use
Taylor's formula at $\theta_i$, where $G_i=0$ and $G_i'=1/\mathsf m_i$. For a
Hermitian matrix apply this to the eigenvalues $u_k$, and
$(\sum_k(u_k-\theta_i)^4)^{1/2}\leq\sum_k(u_k-\theta_i)^2$.
\end{proof}

\begin{lemma}\label{lem:hadamard}
Let $\mathcal M$ be a complete, simply connected Riemannian manifold of
nonpositive sectional curvature, with distance $d$, and $F$ a continuously
differentiable function on $\mathcal M$ that is convex along geodesics. Let
$y(t)$ be a solution of $\dot y=-\operatorname{grad}F(y)$ and $x(t)$ a
continuously differentiable curve with
$\dot x=-\operatorname{grad}F(x)+e(t)$ on an interval $[t_0,t_1]$. Then
$d(x(t),y(t))\leq d(x(t_0),y(t_0))+\int_{t_0}^t|e|$ for $t_0\leq t\leq t_1$.
\end{lemma}

\begin{proof}
On such a manifold any two points $p,p'$ are joined by a unique geodesic
$\gamma_{p,p'}:[0,1]\to\mathcal M$, which depends smoothly on $(p,p')$, and
$d(p,p')^2=\int_0^1|\gamma_{p,p'}'|^2$ is smooth. By the first variation
formula, with $\gamma_t=\gamma_{x(t),y(t)}$ and $\rho(t)=d(x(t),y(t))$,
\[
 \frac d{dt}\,\frac{\rho^2}2=\inner{\dot y}{\gamma_t'(1)}-\inner{\dot x}{\gamma_t'(0)}
 =-\inner{\operatorname{grad}F(y)}{\gamma_t'(1)}+\inner{\operatorname{grad}F(x)}{\gamma_t'(0)}
 -\inner e{\gamma_t'(0)} .
\]
The function $u\mapsto F(\gamma_t(u))$ is convex and differentiable, so its
derivative $\inner{\operatorname{grad}F(\gamma_t(u))}{\gamma_t'(u)}$ is
nondecreasing, and the right side is at most $|e|\,\rho$. For
$\epsilon>0$ this gives $\frac d{dt}\sqrt{\rho^2+\epsilon}\leq|e|$; integrate
and let $\epsilon\to0$.
\end{proof}

For $e=0$ the lemma says that the gradient flow does not increase distances;
Harada and Wilkin prove a related contraction estimate for the trace of the relative metric and its inverse \cite[Theorem~3.18]{hw11}; see also \cite[\S5.3]{hkkp23}.

\begin{theorem}[Comparison with the gradient flow]\label{thm:comparison}
Assume (P) and let $\mathcal H$ be semistable. Let $\mathfrak g(t)$,
$t\geq\tau_0$, be the solution of the weighted gradient flow
\eqref{eq:gradflow} with $\mathfrak g(\tau_0)=\mathfrak h(\tau_0)$. Then for
all $t\geq\tau_0$
\[
 d_{\mathsf m}\bigl(\mathfrak h(t),\mathfrak g(t)\bigr)\leq C_{\mathrm e}
 :=\max_i\frac{3\sqrt3}{16\,\sigma\,y_i^{3/2}}\Bigl(2\Gamma\bigl(\cK_{\mathcal H}(\mathfrak h(\tau_0))-\cK_*\bigr)
 +\frac{\Gamma^2c_\delta^2}{2\kappa}\Bigr)
 +\frac{c_\delta}\kappa\,\max_i\sqrt{\mathsf m_i},
\]
with $\Gamma$ as in \eqref{eq:Gamma}; and, as Hermitian forms,
$e^{-c}\mathfrak g_i(t)\leq\mathfrak h_i(t)\leq e^{c}\mathfrak g_i(t)$ with
$c=C_{\mathrm e}/\min_i\sqrt{\mathsf m_i}$.
\end{theorem}

\begin{proof}
The velocity of $\mathfrak h$ is minus the gradient of $\cK_{\mathcal H}$ for
the metric $\norm\cdot_{\mathsf m}$ plus the tangent vector
$e_i=[G_i(\mu_i)-(\mu_i-\theta_i)/\mathsf m_i]-\delta_i$. By
Lemma~\ref{lem:quadratic},
$\norm e_{\mathsf m}\leq\max_i\frac{3\sqrt3}{16\,\sigma y_i^{3/2}}\norm{\mu-\theta}^2+\max_i\sqrt{\mathsf m_i}\,\norm\delta$.
By Proposition~\ref{prop:perturbed}(i) and Young's inequality,
$\frac d{dt}\cK_{\mathcal H}(\mathfrak h(t))\leq-\frac1{2\Gamma}\norm{\mu-\theta}^2+\frac\Gamma2\norm\delta^2$,
and $\cK_{\mathcal H}\geq\cK_*>-\infty$ because $\mathcal H$ is semistable
(Lemma~\ref{lem:king}(a)); so
$\int_{\tau_0}^\infty\norm{\mu-\theta}^2\leq2\Gamma(\cK_{\mathcal H}(\mathfrak h(\tau_0))-\cK_*)+\Gamma^2c_\delta^2/(2\kappa)$.
Lemma~\ref{lem:hadamard} on $\mathcal X$ with the metric
$\norm\cdot_{\mathsf m}$ and $F=\cK_{\mathcal H}$, which is convex along the
geodesics $u\mapsto\mathfrak he^{u\xi}$ (Lemma~\ref{lem:variation}), gives
the first claim. The second follows from
$\norm{\log(\mathfrak g_i^{-1/2}\mathfrak h_i\mathfrak g_i^{-1/2})}\leq d_{\mathsf m}(\mathfrak h,\mathfrak g)/\sqrt{\mathsf m_i}$.
\end{proof}

\begin{corollary}\label{cor:presentations}
Two curves that satisfy (P) for the same semistable representation
$\mathcal H$ and the same charges remain at bounded distance from each
other. In particular this holds for the normalized metrics induced on the
cohomology of two complexes with a surjective or injective differential at
one vertex whose cohomology representations are identified by a fixed
isomorphism, and for any two initial metrics.
\end{corollary}

\begin{proof}
Each of the two curves is at bounded distance from a solution of
\eqref{eq:gradflow} by Theorem~\ref{thm:comparison}, and two solutions of
\eqref{eq:gradflow} do not move apart by Lemma~\ref{lem:hadamard} with $e=0$.
\end{proof}

The asymptotics of the flow \eqref{eq:gradflow} on a semistable
representation are the subject of \cite{hkkp23}. In the notation of
\cite[\S5.2]{hkkp23} the trace is $\tau(b)=\sum_i\mathsf m_i\tr b_i$, the
central element $\rho$ has the components $\theta_i/\mathsf m_i$, the flow is
\cite[(5.28)]{hkkp23}, and the standing assumption \cite[(5.59)]{hkkp23} is
the semistability of $\mathcal H$ in the sense of
Section~\ref{sec:convergence}. A curve $\mathfrak h$ is called there an
\emph{asymptotic solution} if $C^{-1}\mathfrak g\leq\mathfrak h\leq C\mathfrak g$
for large times, for some, equivalently for every, solution $\mathfrak g$
\cite[\S5.3]{hkkp23}. In \cite[\S5.4]{hkkp23} the lattice of
subrepresentations of $\mathcal H$ of the same phase carries the
\emph{iterated weight filtration}, a filtration labelled by finite sequences
$\lambda=(\lambda_1,\dots,\lambda_n)$ of real numbers and defined
algebraically from that lattice and the function
$\mathsf R\mapsto\sum_i\mathsf m_i\dim\mathsf R_i$
\cite[\S4 and (5.58)]{hkkp23}. After the arrows have been conjugated by a
suitable invertible tuple of matrices, with $1=\sum_\lambda p_\lambda$ the
orthogonal splitting given by the filtration, the function
\begin{equation}\label{eq:iterated}
 \mathfrak a(t)=\sum_{\lambda}t^{\lambda_1}(\log t)^{\lambda_2}\cdots\bigl(\log^{(n-1)}t\bigr)^{\lambda_n}\,p_\lambda
\end{equation}
is an asymptotic solution \cite[Corollary~5.10]{hkkp23}, and every solution
satisfies the bounds of \cite[Theorem~5.11]{hkkp23} on the steps of the
filtration.

\begin{corollary}[Iterated logarithms]\label{cor:iterated}
Assume (P) and let $\mathcal H$ be semistable. Then $\mathfrak h$ is an
asymptotic solution of the gradient flow \cite[(5.28)]{hkkp23} with the
weights $\mathsf m_i=y_i/\sigma^2$. Hence there are an invertible tuple of
matrices $b$ and a constant $C$ such that
$C^{-1}b^*\mathfrak a(t)b\leq\mathfrak h(t)\leq Cb^*\mathfrak a(t)b$ for large
$t$, where $\mathfrak a$ is the function \eqref{eq:iterated} for the arrows
of $\mathcal H$ conjugated by $b$; and the bounds of
\cite[Theorem~5.11]{hkkp23} hold for $\mathfrak h$. The weight function of
the iterated weight filtration is
\[
 \sum_i\mathsf m_i\dim\mathsf R_i=\frac{\Ima Z(\mathsf R)}{\sigma^2}=\frac{|Z(\mathsf R)|}\sigma
\]
on the subquotients $\mathsf R$ of $\mathcal H$ that are semistable of phase
$\varphi$: a constant multiple of the mass.
\end{corollary}

\begin{proof}
The first statement is the last statement of
Theorem~\ref{thm:comparison}; the others follow from the definition of an
asymptotic solution and the results of \cite{hkkp23} just quoted. For a
projection $p$ the inequalities
$e^{-c}\mathfrak g\leq\mathfrak h\leq e^{c}\mathfrak g$ give
$e^{-c}p\mathfrak gp\leq p\mathfrak hp\leq e^{c}p\mathfrak gp$.
\end{proof}

The first statement of the corollary can also be deduced from
\cite[Proposition~5.7]{hkkp23}, which says that a curve satisfying the
equation of the gradient flow up to an absolutely integrable term is an
asymptotic solution; the term is the vector $e$ of the proof of
Theorem~\ref{thm:comparison}. The bounded factor between $\mathfrak h$ and
\eqref{eq:iterated} is not shown to converge, also not in \cite{hkkp23} for
the gradient flow; it converges in Proposition~\ref{prop:weights} below and
in the example of \cite[Proposition~6.9]{wue26b}.
In the following example the exponents determine the ratio of the vertex
weights.

\begin{proposition}\label{prop:weights}
Let $Q$ be the quiver $0\to1$, let $z_0,z_1$ have the same argument
$\varphi$, and let $\mathcal H=(\CC\xrightarrow{b}\CC)$ with $b\neq0$; it is
semistable and not polystable, a nonsplit extension of the simple
representation at the vertex $0$ by the simple representation at the vertex
$1$. For every initial metric the normalized phase flow
$\mathfrak h=e^{\varphi t}h$ satisfies
\[
 \log\mathfrak h_0(t)=\frac{|z_1|}{|z_0|+|z_1|}\log t+c_0+o(1),\qquad
 \log\mathfrak h_1(t)=-\frac{|z_0|}{|z_0|+|z_1|}\log t+c_1+o(1)
\]
as $t\to\infty$, with constants $c_0,c_1$. The solutions of
\eqref{eq:gradflow} with arbitrary positive weights $\mathsf m_0',\mathsf m_1'$
have the exponents $\mathsf m_1'/(\mathsf m_0'+\mathsf m_1')$ and
$-\mathsf m_0'/(\mathsf m_0'+\mathsf m_1')$.
\end{proposition}

\begin{proof}
Here $\theta=0$ and $\mathsf m_i=|z_i|/\sigma$. The squared norm
$q=|b|^2\mathfrak h_1/\mathfrak h_0$ of the arrow gives $\mu_0=q$,
$\mu_1=-q$, so $(\log\mathfrak h_0)'=G_0(q)$, $(\log\mathfrak h_1)'=G_1(-q)$
and $(\log q)'=-(G_0(q)-G_1(-q))<0$. By Lemma~\ref{lem:quadratic},
$G_0(q)-G_1(-q)=\varkappa q+O(q^2)$ with
$\varkappa=1/\mathsf m_0+1/\mathsf m_1$. So $q$ decreases to zero,
$\frac d{dt}(1/q)=\varkappa+O(q)$, hence $1/q=\varkappa t+O(\log t)$ and
$q=1/(\varkappa t)+O(t^{-2}\log t)$. Then
$(\log\mathfrak h_0)'=q/\mathsf m_0+O(q^2)$ integrates to
$\frac1{\mathsf m_0\varkappa}\log t+c_0+o(1)$, and
$\frac1{\mathsf m_0\varkappa}=\frac{\mathsf m_1}{\mathsf m_0+\mathsf m_1}$;
likewise for $\mathfrak h_1$. For \eqref{eq:gradflow} with weights
$\mathsf m_i'$ the same computation applies with
$(\log\mathfrak g_0)'=q/\mathsf m_0'$, $(\log\mathfrak g_1)'=-q/\mathsf m_1'$.
\end{proof}

So on this representation the metric grows like $t$ to the power
$|Z(\mathsf S_1)|/|Z(\mathcal H)|$ on the quotient $\mathsf S_0$ and decays
like $t$ to the power $-|Z(\mathsf S_0)|/|Z(\mathcal H)|$ on the
subrepresentation $\mathsf S_1$, where $\mathsf S_i$ is the simple
representation at the vertex $i$; and a gradient flow whose weights are not
proportional to the masses $|z_i|$ is at logarithmically growing distance
from the phase flow. For a generalized Kronecker quiver with
$\alpha_0\neq\alpha_1$ the weights do not matter, because the stable
factors of a semistable representation then have proportional dimension
vectors.

\subsection{The mass and the Harder--Narasimhan filtration}
A representation $\mathsf R\neq0$ is semistable of phase
$\phi=\arg Z(\mathsf R)$, in the sense of Section~\ref{sec:convergence} with
the parameter $\theta$ formed with $\phi$, if and only if
$\arg Z(\mathsf R')\leq\phi$ for all nonzero subrepresentations
$\mathsf R'$. Every nonzero representation $\mathcal H$ has a unique
filtration
\[
 0=\mathcal H^0\subset\mathcal H^1\subset\dots\subset\mathcal H^{\mathsf n}=\mathcal H
\]
by subrepresentations whose factors
$\operatorname{gr}^k=\mathcal H^k/\mathcal H^{k-1}$ are semistable of phases
$\phi_1>\phi_2>\dots>\phi_{\mathsf n}$, the \emph{Harder--Narasimhan
filtration} for the stability condition of Section~\ref{sec:setting}
\cite[Proposition~2.4]{bri07}; a filtration by subrepresentations with
semistable factors of strictly decreasing phases is this one, because there
are no nonzero morphisms from a semistable representation to a semistable
representation of smaller phase. We call
$\mass_{\mathrm{HN}}(\mathcal H)=\sum_k|Z(\operatorname{gr}^k)|$ the
\emph{mass of the filtration}; it is the mass of the object $\mathcal H$ for
the stability condition. For a tuple of metrics on $\mathcal H$ let
$\mass(\mathfrak h)=\sum_i\tr|z_i+\mu_i|$ be the mass of the metrized
representation, $\Pi_k$ the orthogonal projections onto $\mathcal H^k$, and
\[
 \mathcal Y_k=\sum_a\norm{\Pi_{k,\mathsf t(a)}\beta_a(1-\Pi_{k,\mathsf s(a)})}^2
\]
the squared norm of the second fundamental form of $\mathcal H^k$.

\begin{proposition}[Dissipation of the mass]\label{prop:dissipation}
Assume (P). Then for $t\geq\tau_0$
\[
 \frac{d\,\mass(\mathfrak h(t))}{dt}\leq-\mathfrak D+4\bar\cN\norm\delta,
 \qquad
 \mathfrak D\geq\min(\sin\psi_-,\sin\psi_+)\sum_a\norm{\beta_a\Phi^{\mathcal H}_{\mathsf s(a)}-\Phi^{\mathcal H}_{\mathsf t(a)}\beta_a}^2,
\]
where
$\mathfrak D=\sum_a\sum_{j,l}|\inner{v_l}{\beta_au_j}|^2(p_j-q_l)(\cos q_l-\cos p_j)$
for orthonormal eigenvectors $u_j$ of $\Phi^{\mathcal H}_{\mathsf s(a)}$ and
$v_l$ of $\Phi^{\mathcal H}_{\mathsf t(a)}$ with eigenvalues $p_j$, $q_l$. In
particular the mass is nonincreasing along the phase flow of a
representation, for all charges.
\end{proposition}

\begin{proof}
The derivative of $u\mapsto|z_i+u|$ is $\cos\arg(z_i+u)$, so
$\frac d{dt}\mass=\sum_i\tr(\cos\Phi^{\mathcal H}_i\,\dot\mu_i)$. Insert
\eqref{eq:betaflow}; the result is linear in
$\Phi^{\mathcal H}+\delta$. The part with $\Phi^{\mathcal H}$ is, since
$\Phi^{\mathcal H}_i$ commutes with its cosine,
\[
 \sum_a\sum_{j,l}|\inner{v_l}{\beta_au_j}|^2\bigl(p_j\cos p_j-q_l\cos p_j-p_j\cos q_l+q_l\cos q_l\bigr)=-\mathfrak D .
\]
The cosine decreases on $(0,\pi)$, and
$(p-q)(\cos q-\cos p)\geq\min(\sin\psi_-,\sin\psi_+)(p-q)^2$ for
$p,q\in[\psi_-,\psi_+]$. The part with $\delta$ consists of four traces for
each arrow, each at most $\norm{\beta_a}^2\norm\delta$.
\end{proof}

For Kronecker quivers the monotonicity of the mass on representations is
\cite[Proposition~9.1]{wue26b}.

\begin{proposition}[Mass of a metric and mass of the filtration]\label{prop:hnmass}
For every tuple of metrics on $\mathcal H$,
\[
 \mass(\mathfrak h)\geq\mass_{\mathrm{HN}}(\mathcal H)
 +\sum_{k=1}^{\mathsf n-1}\bigl(\cos\phi_{k+1}-\cos\phi_k\bigr)\,\mathcal Y_k ,
\]
and the coefficients $\cos\phi_{k+1}-\cos\phi_k$ are positive.
\end{proposition}

\begin{proof}
Let $Q_k=\Pi_k-\Pi_{k-1}$. By the pinching inequality for the trace norm and
$\tr|X|\geq|\tr X|$,
$\mass\geq\sum_k\bigl|\sum_i\tr(Q_{k,i}(z_i+\mu_i))\bigr|$. For a
subrepresentation with projections $\Pi$ one has
$\sum_i\tr(\Pi_i\mu_i)=\sum_a(\norm{\beta_a\Pi_{\mathsf s(a)}}^2-\norm{\Pi_{\mathsf t(a)}\beta_a}^2)
=-\sum_a\norm{\Pi_{\mathsf t(a)}\beta_a(1-\Pi_{\mathsf s(a)})}^2$, as in the
proof of Lemma~\ref{lem:king}(b). So with
$\mathcal Y_0=\mathcal Y_{\mathsf n}=0$,
$\sum_i\tr(Q_{k,i}(z_i+\mu_i))=Z(\operatorname{gr}^k)-(\mathcal Y_k-\mathcal Y_{k-1})$,
and the modulus of this number is at least its component in the direction of
$Z(\operatorname{gr}^k)$, which is
$|Z(\operatorname{gr}^k)|-(\mathcal Y_k-\mathcal Y_{k-1})\cos\phi_k$. Sum
over $k$ and sum by parts.
\end{proof}

\subsection{Unstable cohomology}
Assume (P). The arrows $\beta(t)=(\beta_a(t))$ in the frames of
\eqref{eq:betaflow} form a bounded curve. Let $\omega$ be the set of its
limit points as $t\to\infty$; it is nonempty, compact and connected.

\begin{lemma}\label{lem:limitset}
Assume \emph{(P)}.
\begin{enumerate}
\item[(a)] The mass $\mass(\mathfrak h(t))$ converges, and
 $\beta_a\Phi^{\mathcal H}_{\mathsf s(a)}-\Phi^{\mathcal H}_{\mathsf t(a)}\beta_a$,
 $\dot\beta_a$ and the derivatives of the $\Phi^{\mathcal H}_i$ tend to
 zero.
\item[(b)] There are numbers $\pi>\phi_1'>\dots>\phi_{\mathsf n'}'>0$ and
 nonzero dimension vectors $d_k=(n_{i,k})_i$ with $\sum_kd_k=(n_i)$ and
 $\arg Z(d_k)=\phi_k'$ such that for every $\beta^\infty\in\omega$ the phase
 of $\beta^\infty$ at the vertex $i$ has the eigenvalues $\phi_k'$ with
 multiplicities $n_{i,k}$, and the representation with the arrows
 $\beta^\infty$ is the orthogonal direct sum of representations
 $\mathsf P_k$ of dimension vectors $d_k$, supported on the eigenspaces,
 with the moment map of a harmonic metric of phase $\phi_k'$ on
 $\mathsf P_k$. In particular $\mathsf P_k$ is polystable of phase
 $\phi_k'$.
\item[(c)] The eigenvalues of $\Phi^{\mathcal H}_i(t)$, with multiplicity,
 converge to the $\phi_k'$ with multiplicities $n_{i,k}$, and
 $\mass(\mathfrak h(t))\to\sum_k|Z(d_k)|$.
\end{enumerate}
\end{lemma}

\begin{proof}
(a) By Proposition~\ref{prop:dissipation},
$\mass+4\bar\cN\int_t^\infty\norm\delta$ is nonincreasing and nonnegative,
so $\mass$ converges and $\mathfrak D$ is integrable. By
\eqref{eq:betaflow} the $\dot\beta_a$ are bounded, hence so are the
derivatives of the $\mu_i$ and, by the Lipschitz property, of the
$\Phi^{\mathcal H}_i$. So
$\sum_a\norm{\beta_a\Phi^{\mathcal H}_{\mathsf s(a)}-\Phi^{\mathcal H}_{\mathsf t(a)}\beta_a}^2$
is Lipschitz and integrable and tends to zero; then $\dot\beta_a\to0$ by
\eqref{eq:betaflow}, and the derivatives of $\mu_i$ and
$\Phi^{\mathcal H}_i$ tend to zero.

(b) For $\beta^\infty\in\omega$, by (a) and continuity each arrow
intertwines the phases at its source and target. So the eigenspaces of the
phase for a fixed eigenvalue $p$ form an orthogonal direct summand
$\mathsf P_p$, and on it $\arg(z_i+\mu_i)=p$ for all $i$: the metric is
harmonic of phase $p$ on $\mathsf P_p$, which is therefore polystable
(Lemma~\ref{lem:king}(b)), and $\arg Z(\mathsf P_p)=p$ because
$\sum_i\tr(\Pi_i\mu_i)=0$ for an orthogonal direct summand. Thus the
eigenvalues of the phase of $\beta^\infty$ lie in the finite set of
arguments of the numbers $Z(d)$, $0\neq d\leq(n_i)$. The ordered list of
these eigenvalues at each vertex depends continuously on $\beta$, takes
values in a finite set on $\omega$, and $\omega$ is connected; so it is
constant on $\omega$.

(c) The distance from $\beta(t)$ to $\omega$ tends to zero, and on
$\mathsf P_k$ all eigenvalues of $z_i+\mu_i$ lie on one ray.
\end{proof}

By (c) there is a time after which the spectrum of each
$\Phi^{\mathcal H}_i(t)$ consists of separated clusters near the
$\phi_k'$. Let $P_k(t)$ be the spectral projection of the $k$-th cluster and
$P_{\leq k}=P_1+\dots+P_k$. Since the derivatives of the
$\Phi^{\mathcal H}_i$ tend to zero and the clusters stay separated,
\begin{equation}\label{eq:slow}
 \norm{\dot P_{\leq k}(t)}\to0,\qquad\norm{P_l\beta_aP_k}\to0\quad(l\neq k),
\end{equation}
the second because $X=P_l\beta_aP_k$ satisfies
$X\Phi^{\mathcal H}_{\mathsf s(a)}-\Phi^{\mathcal H}_{\mathsf t(a)}X=P_l(\beta_a\Phi^{\mathcal H}_{\mathsf s(a)}-\Phi^{\mathcal H}_{\mathsf t(a)}\beta_a)P_k$
with separated spectra on the two sides.

\begin{lemma}\label{lem:dichotomy}
Let $\Psi(t)$, $t\geq T$, be a continuous family of Hermitian
$n\times n$ matrices and $P(t)$ a continuously differentiable family of
orthogonal projections of constant rank. Suppose there are numbers
$a_+>a_-$ and $\epsilon\geq0$ with $12\epsilon\leq a_+-a_-$ such that for
all $t\geq T$ and all vectors $w$, with $w'=P(t)w$ and $w''=w-w'$,
\[
 \Rea\inner{w'}{\Psi w}\geq a_+|w'|^2-\epsilon|w'||w''|,\qquad
 \Rea\inner{w''}{\Psi w}\leq a_-|w''|^2+\epsilon|w'||w''|,\qquad
 \norm{\dot P}\leq\epsilon .
\]
Then $\CC^n=S^{\mathrm f}\oplus S^{\mathrm s}$ with
$\dim S^{\mathrm f}=\operatorname{rank}P$, such that for the solutions of
$\dot w=-\frac12\Psi w$ on $[T,\infty)$:
\begin{enumerate}
\item[(i)] if $w(T)\in S^{\mathrm f}$, then
 $|w''(t)|\leq\frac{6\epsilon}{a_+-a_-}|w'(t)|$ and
 $|w(t)|^2\leq2e^{-(a_+-3\epsilon)(t-T)}|w(T)|^2$ for $t\geq T$;
\item[(ii)] if $w(T)\in S^{\mathrm s}$, then
 $|w(t)|^2\geq\frac12e^{-(a_-+3\epsilon)(t-T)}|w(T)|^2$ for $t\geq T$.
\end{enumerate}
\end{lemma}

\begin{proof}
Since $P\dot PP=0=(1-P)\dot P(1-P)$, a solution satisfies
$\frac d{dt}|w'|^2\leq-a_+|w'|^2+3\epsilon|w'||w''|$ and
$\frac d{dt}|w''|^2\geq-a_-|w''|^2-3\epsilon|w'||w''|$. Let
$c_0=6\epsilon/(a_+-a_-)\leq\frac12$. For $c\in[c_0,1]$ with $c>0$ and a
nonzero solution, at a time with $|w''|=c|w'|$,
\[
 \frac d{dt}\bigl(c^2|w'|^2-|w''|^2\bigr)\leq c^2|w'|^2\bigl(-(a_+-a_-)+3\epsilon(c+c^{-1})\bigr)<0,
\]
because $3\epsilon c\leq\frac14(a_+-a_-)$ and
$3\epsilon c^{-1}\leq\frac12(a_+-a_-)$. So a solution in the open cone
$\{|w''|>c|w'|\}$ cannot reach its boundary: these cones are invariant
forwards in time. Let $S^{\mathrm s}=\ker P(T)$. A solution starting there
stays in the cone with $c=1$, where
$\frac d{dt}|w''|^2\geq-(a_-+3\epsilon)|w''|^2$; this gives (ii). For
$\tau\geq T$ let $S_\tau$ be the space of the initial values $w(T)$ of the
solutions with $w(\tau)$ in the range of $P(\tau)$. Such a solution
satisfies $|w''|\leq c_0|w'|$ on $[T,\tau]$, since otherwise it would lie, at
the time $\tau$, in one of the open cones with $c\in[c_0,1]$, $c>0$. Let $S^{\mathrm f}$ be a
limit point of the spaces $S_\tau$ in the Grassmannian as $\tau\to\infty$.
By continuity every solution starting in $S^{\mathrm f}$ satisfies
$|w''|\leq c_0|w'|$ for all $t\geq T$, and then
$\frac d{dt}|w'|^2\leq-(a_+-3\epsilon)|w'|^2$ and $|w|^2\leq2|w'|^2$; this is
(i). The two growth estimates exclude a common nonzero vector, and the
dimensions add up to $n$.
\end{proof}

\begin{theorem}[Harder--Narasimhan asymptotics]\label{thm:hn}
Let $\mathcal H$ be a nonzero representation of $Q$, with
Harder--Narasimhan filtration $(\mathcal H^k)$, factors
$\operatorname{gr}^k$ and phases $\phi_1>\dots>\phi_{\mathsf n}$, and let
$\mathfrak h(t)$ satisfy \emph{(P)}. Then:
\begin{enumerate}
\item[(i)] the eigenvalues of $\Phi^{\mathcal H}_i(t)$, with multiplicity,
 converge to the $\phi_k$ with multiplicities
 $\dim\operatorname{gr}^k_i$, and $\mass(\mathfrak h(t))$ converges to
 $\mass_{\mathrm{HN}}(\mathcal H)$, the infimum of the mass over all metrics
 on $\mathcal H$;
\item[(ii)] for every nonzero $v\in\mathcal H_i$,
 \[
 \lim_{t\to\infty}\frac1t\log\mathfrak h_i(t)(v,v)=\varphi-\phi_k,
 \]
 where $k$ is the least index with $v\in\mathcal H^k_i$; so the
 Harder--Narasimhan filtration is the filtration of the spaces by the
 exponential rate of decay of the metric;
\item[(iii)] for every $\epsilon>0$ there is a constant $C$ with
 $\mathcal Y_k(t)\leq Ce^{-2\varkappa_kt}$,
 $\varkappa_k=\min\{\frac12(\phi_k-\phi_{k+1})-\epsilon,\ \kappa\}$, for
 $k<\mathsf n$ (with $\kappa=\infty$ if $\delta=0$);
\item[(iv)] for each $k$ the metric induced by $\mathfrak h(t)$ on
 $\operatorname{gr}^k$, multiplied by $e^{(\phi_k-\varphi)t}$, satisfies
 \emph{(P)} for the semistable representation $\operatorname{gr}^k$. For one
 factor the original perturbation and its rate $\kappa$ are retained.
 For two or more factors take $0<2\epsilon<\min_{l<\mathsf n}(\phi_l-\phi_{l+1})$;
 the perturbation then has rate $\min_{l<\mathsf n}\varkappa_l$. Hence
 Proposition~\ref{prop:perturbed}, Theorem~\ref{thm:comparison} and
 Corollary~\ref{cor:iterated} apply to it: its moment map tends to the
 stability parameter of $\operatorname{gr}^k$, it is an asymptotic solution
 of the weighted gradient flow of $\operatorname{gr}^k$, and if
 $\operatorname{gr}^k$ is polystable it converges exponentially to a
 harmonic metric of $\operatorname{gr}^k$.
\end{enumerate}
\end{theorem}

\begin{proof}
\emph{Step 1: the filtration by exponents.} Fix $k<\mathsf n'$, a vertex
$i$, and $\epsilon>0$ with $60\epsilon\leq\phi_k'-\phi_{k+1}'$. By
Lemma~\ref{lem:limitset} and \eqref{eq:slow} there is a time $T$ after which
the clusters are within $\epsilon$ of the $\phi_l'$,
$\norm\delta\leq\epsilon$ and $\norm{\dot P_{\leq k}}\leq\epsilon$. For a
vector $v\in\mathcal H_i$ the coordinates $w$ of $v$ in the frame satisfy
$|w|^2=\mathfrak h_i(v,v)$ and $\dot w=-\frac12\Psi w$ with
$\Psi=\Phi^{\mathcal H}_i-\varphi+\delta_i$. Since $P=P_{\leq k}$ commutes
with $\Phi^{\mathcal H}_i$, the hypotheses of Lemma~\ref{lem:dichotomy} hold
with $a_+=\phi_k'-\varphi-2\epsilon$ and
$a_-=\phi_{k+1}'-\varphi+2\epsilon$. In terms of vectors of $\mathcal H_i$
the lemma gives complementary subspaces, of which the first has the
dimension $\sum_{l\leq k}n_{i,l}$, with
$\mathfrak h_i(t)(v,v)\leq C_ve^{-(\phi_k'-\varphi-5\epsilon)t}$ on the
first and
$\mathfrak h_i(t)(v,v)\geq c_ve^{-(\phi_{k+1}'-\varphi+5\epsilon)t}$,
$c_v>0$, for the nonzero vectors of the second. A vector with a nonzero
component in the second subspace then satisfies
$\liminf\frac1t\log\mathfrak h_i(t)(v,v)\geq-(\phi_{k+1}'-\varphi+5\epsilon)$.
So the first subspace is the set of all $v$ with
$\limsup\frac1t\log\mathfrak h_i(t)(v,v)\leq-(\phi_k'-\varphi-5\epsilon)$;
these sets decrease with $\epsilon$ and have the same dimension, so they
coincide with
\[
 F_{k,i}:=\Bigl\{v\in\mathcal H_i:\ \limsup_{t\to\infty}\frac1t\log\mathfrak h_i(t)(v,v)\leq\varphi-\phi_k'\Bigr\},
 \qquad\dim F_{k,i}=\sum_{l\leq k}n_{i,l},
\]
and for $v\in F_{k,i}\setminus F_{k-1,i}$ the limit of
$\frac1t\log\mathfrak h_i(t)(v,v)$ exists and equals $\varphi-\phi_k'$
(with $F_0=0$ and $F_{\mathsf n'}=\mathcal H$). Here the lower bound for
$k=1$ and the upper bound for $k=\mathsf n'$, which do not come from a gap,
follow from $\frac d{dt}|w|^2=-\inner w{\Psi w}$ and
Lemma~\ref{lem:limitset}(c): the eigenvalues of $\Psi$ lie eventually between
$\phi_{\mathsf n'}'-\varphi-2\epsilon$ and $\phi_1'-\varphi+2\epsilon$.

\emph{Step 2: the $F_k$ are subrepresentations.} For an arrow $a$ of
$\mathcal H$ and $v\in\mathcal H_{\mathsf s(a)}$, the norm of the image of
$v$ for $\mathfrak h_{\mathsf t(a)}$ is at most $\bar\cN^{1/2}$ times the
norm of $v$ for $\mathfrak h_{\mathsf s(a)}$. So the arrows do not decrease
the rate of decay, and they map $F_k$ to $F_k$. The factor $F_k/F_{k-1}$ has
the dimension vector $d_k$, hence $\arg Z(F_k/F_{k-1})=\phi_k'$.

\emph{Step 3.} By Lemma~\ref{lem:dichotomy}(i), applied with $\epsilon$
arbitrarily small for large $T$, the coordinates of the vectors of $F_k$
form the graph of a linear map from the range of $P_{\leq k}(t)$ to its
orthogonal complement whose norm tends to zero. So for the
$\mathfrak h(t)$-orthogonal projections $\Pi_k'(t)$ onto $F_k$, in the
frames, $\norm{\Pi_k'-P_{\leq k}}\to0$; and with
$Q_k=\Pi_k'-\Pi_{k-1}'$, by \eqref{eq:slow} and Lemma~\ref{lem:limitset}(c),
\begin{equation}\label{eq:blockslimit}
 Q_l\beta_aQ_k\to0\ (l\neq k),\qquad
 \norm{Q_k(\Phi^{\mathcal H}-\phi_k')Q_k}\to0,\qquad
 Q_l\Phi^{\mathcal H}Q_k\to0\ (l\neq k).
\end{equation}

\emph{Step 4: the factors are semistable.} In frames adapted to the
decomposition $1=\sum_kQ_k(t)$ the arrows are block upper triangular, and the
$k$-th diagonal block is the tuple of arrows of the representation
$F_k/F_{k-1}$ in orthonormal frames of its induced metric, a point of the
orbit of this representation under the changes of bases. Along a suitable
sequence of times these frames and the arrows converge, and by
\eqref{eq:blockslimit} the $k$-th diagonal block converges to the arrows of
the summand $\mathsf P_k$ of a point of $\omega$. So the closure of the orbit
of $F_k/F_{k-1}$ contains a representation isomorphic to $\mathsf P_k$, which
is semistable of phase $\phi_k'$. If $F_k/F_{k-1}$ had a subrepresentation
of a dimension vector $d'$ with $\arg Z(d')>\phi_k'$, then so would every
representation in the closure of its orbit, because the set of
representations with a subrepresentation of dimension vector $d'$ is the
image of a closed set under a projection with compact fibres, a product of
Grassmannians. This contradicts the semistability of $\mathsf P_k$. So
$F_k/F_{k-1}$ is semistable of phase $\phi_k'$.

\emph{Step 5.} The filtration $(F_k)$ has semistable factors of strictly
decreasing phases. Hence it is the Harder--Narasimhan filtration:
$\mathsf n'=\mathsf n$, $F_k=\mathcal H^k$, $\phi_k'=\phi_k$ and
$d_k=\dim\operatorname{gr}^k$. With Lemma~\ref{lem:limitset}(c) and Step~1
this gives (i) and (ii); the mass of the filtration is the infimum of the
mass by Proposition~\ref{prop:hnmass}.

\emph{Step 6: the second fundamental forms.} Fix $k<\mathsf n$ and write
$\Pi=\Pi_k$. For an arrow $a$ let
$\upsilon_a=\Pi_{\mathsf t(a)}\beta_a(1-\Pi_{\mathsf s(a)})$, and for a
Hermitian matrix $M$ at a vertex let $M^{\mathrm{off}}=\Pi M(1-\Pi)+(1-\Pi)M\Pi$.
With $\tilde\Phi=\Phi^{\mathcal H}+\delta$ we claim that
\begin{equation}\label{eq:secondform}
 \frac{d\mathcal Y_k}{dt}=\sum_a\Bigl[\tr\bigl(\upsilon_a^\dagger\upsilon_a\,\tilde\Phi_{\mathsf s(a)}\bigr)
 -\tr\bigl(\upsilon_a\upsilon_a^\dagger\,\tilde\Phi_{\mathsf t(a)}\bigr)\Bigr]
 +\sum_i\tr\bigl(\tilde\Phi_i^{\mathrm{off}}\mu_i^{\mathrm{off}}\bigr).
\end{equation}
Indeed, the coordinate maps of the frames satisfy
$\dot gg^{-1}=-\frac12(\tilde\Phi-\varphi)$, a Hermitian matrix, and the
range of $\Pi$ is the image of a fixed subspace, so
$\dot\Pi=-\frac12\tilde\Phi^{\mathrm{off}}$. Differentiate
$\upsilon_a$ with \eqref{eq:betaflow}, use
$(1-\Pi_{\mathsf t(a)})\beta_a\Pi_{\mathsf s(a)}=0$, and take the trace
against $\upsilon_a^\dagger$: with
$\beta_a'=\Pi_{\mathsf t(a)}\beta_a\Pi_{\mathsf s(a)}$ and
$\beta_a''=(1-\Pi_{\mathsf t(a)})\beta_a(1-\Pi_{\mathsf s(a)})$,
\begin{align*}
 \frac d{dt}\norm{\upsilon_a}^2={}&\tr\bigl(\upsilon_a^\dagger\upsilon_a\tilde\Phi_{\mathsf s(a)}\bigr)
 -\tr\bigl(\upsilon_a\upsilon_a^\dagger\tilde\Phi_{\mathsf t(a)}\bigr)\\
 &-2\Rea\tr\bigl(\Pi\tilde\Phi_{\mathsf t(a)}(1-\Pi)\beta_a''\upsilon_a^\dagger\bigr)
 +2\Rea\tr\bigl(\Pi\tilde\Phi_{\mathsf s(a)}(1-\Pi)\upsilon_a^\dagger\beta_a'\bigr),
\end{align*}
and
$(1-\Pi)\mu_i\Pi=\sum_{\mathsf s(a)=i}\upsilon_a^\dagger\beta_a'-\sum_{\mathsf t(a)=i}\beta_a''\upsilon_a^\dagger$;
summing over the arrows gives \eqref{eq:secondform}.

The term $\sum_i\tr((\Phi^{\mathcal H}_i)^{\mathrm{off}}\mu_i^{\mathrm{off}})$
is nonpositive: the matrix $2\Pi-1$ is unitary and Hermitian, conjugation
by it changes the sign of the off-diagonal parts, and
$\Phi^{\mathcal H}_i$ is a decreasing function of $\mu_i$, so four times the
trace is $\tr((g(A)-g(B))(A-B))\leq0$ with $A=\mu_i$,
$B=(2\Pi-1)\mu_i(2\Pi-1)$ and $g=\arg(z_i+\,\cdot\,)$. By
\eqref{eq:blockslimit} and Step~5, for every $\epsilon>0$ and large $t$
the eigenvalues of $\Pi\Phi^{\mathcal H}\Pi$ on the range of $\Pi$ are at
least $\phi_k-\epsilon$ and those of $(1-\Pi)\Phi^{\mathcal H}(1-\Pi)$ on the
complement are at most $\phi_{k+1}+\epsilon$. The terms with $\delta$ are at
most $2\norm\delta\mathcal Y_k+2\sqrt2\,\norm\delta(\bar\cN\mathcal Y_k)^{1/2}$.
Hence
\[
 \frac{d\mathcal Y_k}{dt}\leq-\bigl(\phi_k-\phi_{k+1}-2\epsilon-2\norm\delta\bigr)\mathcal Y_k
 +2\sqrt2\,\norm\delta\,(\bar\cN\mathcal Y_k)^{1/2}
\]
for large $t$. With $\norm\delta\leq c_\delta e^{-\kappa(t-\tau_0)}$ this
gives (iii); one applies the inequality with $\epsilon/2$ in the place of
$\epsilon$, so that a factor $t$ arising when the two exponents coincide is
absorbed.

\emph{Step 7: the factors.} If $\mathsf n=1$, the factor is $\mathcal H$
and $\phi_1=\varphi$, so (iv) follows directly from (P). For $\mathsf n>1$,
the restriction of a metric to a subspace and
the quotient metric move with the compressed velocity. So in the adapted
frames of Step~4 the metric induced on $\operatorname{gr}^k$, multiplied by
$e^{(\phi_k-\varphi)t}$, has the velocity
$\phi_k-Q_k\Phi^{\mathcal H}Q_k-Q_k\delta Q_k$, and its arrows are the $k$-th
diagonal block. Let $\mu^{(k)}$ be the moment map of this block, $D$ the
block-diagonal part of $\mu$ and $O=\mu-D$. Every block of $O$, and every
block of $Q_kDQ_k-\mu^{(k)}$, is a sum of products of two blocks of the
arrows of which at least one, respectively both, are off-diagonal; so with
$Y^2=\sum_a\sum_{l<k}\norm{Q_l\beta_aQ_k}^2\leq\sum_k\mathcal Y_k$,
$\norm O\leq4\bar\cN^{1/2}Y$ and $\norm{Q_kDQ_k-\mu^{(k)}}\leq Y^2$. By the
Lipschitz property, and because a function of $D$ is block diagonal,
\[
 \norm{Q_k\Phi^{\mathcal H}Q_k-\arg(z+\mu^{(k)})}\leq\frac{4\bar\cN^{1/2}Y+Y^2}{\min_iy_i}.
\]
So the velocity is $\phi_k-\arg(z+\mu^{(k)})-\delta^{(k)}$ with
$\delta^{(k)}$ exponentially small by (iii), and the squared norms of the
arrows of the block sum to at most $\bar\cN$: this is (P) for
$\operatorname{gr}^k$, which is semistable of phase $\phi_k$.
\end{proof}

\begin{remark}\label{rem:hn}
(1) For a subrepresentation of a representation, and $\delta=0$, the
identity \eqref{eq:secondform} holds at every time, with $\mathcal Y$ the
squared norm of its second fundamental form. Its last term is nonpositive.
So the second fundamental form decays as long as the phases on the
subrepresentation exceed those on its orthogonal complement.

(2) The filtration by the rates of decay is a filtration by
subrepresentations for the simple reason of Step~2; the theorem identifies
the rates and the dimensions. For the gradient flow the limiting filtration
is described in the same terms in \cite[\S1.1]{hkkp23}, and for unstable
points of Kempf--Ness gradient flows the term linear in $t$ is in
\cite{lam26}.

(3) By (i) and Proposition~\ref{prop:dissipation}, along the phase flow of a
representation of a quiver without oriented cycles the mass decreases to the
mass of the object for the stability condition; the largest and the
smallest eigenvalue of the phase tend to $\phi_1$ and $\phi_{\mathsf n}$.
\end{remark}

\section{Kronecker quivers: complexes with differentials at one vertex}\label{sec:kronecker}

Let $Q=K_m$, with vertices $0,1$ and arrows $a_1,\dots,a_m$ from $0$ to $1$.
A bounded complex of representations $\mathsf E$ has spaces $V^n$ at the
vertex $0$ and $W^n$ at the vertex $1$, differentials
$d_0:V^n\to V^{n+1}$ and $d_1:W^n\to W^{n+1}$, and arrows
$A_j^n:V^n\to W^n$ with $d_1A_j^n=A_j^{n+1}d_0$. Its cohomology in degree $n$
is a representation $H^n(\mathsf E)$. The endomorphism $\Omega$, the phase
and the phase flow of $\mathsf E$ with a metric are those of
Section~\ref{sec:setting}, and we use the properties (O1) and (O2) stated
there. In this section the charges $z_0,z_1$ are arbitrary in the upper
half-plane unless stated otherwise; we call them \emph{strictly ordered} if
$\alpha_0>\alpha_1$, that is, if the simple representation at the vertex $0$
has the larger phase. For strictly ordered charges there are stable
representations with both spaces nonzero; for $\alpha_0\leq\alpha_1$ the only
stable representations are the two simple ones, because a representation
$(\mathsf R_0\to\mathsf R_1)$ with $\mathsf R_0,\mathsf R_1\neq0$ has the subrepresentation $(0\to\mathsf R_1)$, whose phase $\alpha_1$ is then
not smaller than that of the representation.

\subsection{The decomposition}

\begin{theorem}[Decomposition]\label{thm:decomposition}
Let $\mathsf E$ be a bounded complex of representations of $K_m$ with a
metric.
\begin{enumerate}
\item[(a)] Suppose $d_1=0$. Put $U^n=d_0(V^{n-1})\subset V^n$, and let
 $(U^n)^\perp$ be its orthogonal complement in $V^n$. For each $n$ the spaces
 \[
 \bigl((U^n)^\perp\xrightarrow{A_j^n}W^n\bigr)\ \text{in degree }n,\qquad
 (U^{n+1}\to0)\ \text{in degree }n+1,
 \]
 with the differential $d_0:(U^n)^\perp\to U^{n+1}$, which is surjective,
 form a subcomplex of representations $\mathsf E_n$. The complex $\mathsf E$
 is the orthogonal direct sum of the $\mathsf E_n$. If $H^n(\mathsf E)\neq0$,
 then $\mathsf E_n[n]$ is a complex with a surjective differential at the
 vertex $0$ (Definition~\ref{def:type}), with cohomology $H^n(\mathsf E)$.
\item[(b)] Suppose $d_0=0$. Let $C^n$ be the orthogonal complement of
 $\ker d_1^n$ in $W^n$. For each $n$ the spaces
 \[
 (0\to C^{n-1})\ \text{in degree }n-1,\qquad
 \bigl(V^n\xrightarrow{A_j^n}\ker d_1^n\bigr)\ \text{in degree }n,
 \]
 with the differential $d_1:C^{n-1}\to\ker d_1^n$, which is injective, form
 a subcomplex of representations $\mathsf E_n$. The complex $\mathsf E$ is the
 orthogonal direct sum of the $\mathsf E_n$. If $H^n(\mathsf E)\neq0$, then
 $\mathsf E_n[n]$ is a complex with an injective differential at the vertex
 $1$ (Definition~\ref{def:type}), with cohomology $H^n(\mathsf E)$.
\item[(c)] In both cases, if $H^n(\mathsf E)=0$, then $\mathsf E_n$ is a
 complex $C\to C'$ at one vertex with invertible differential, on which the
 arrows vanish.
\item[(d)] In both cases $\Omega$ and the phase of $\mathsf E$ are the direct
 sums of those of the $\mathsf E_n$, and the phase flow preserves the
 decomposition: the subspaces remain orthogonal for the metric at every time,
 and the metric of each $\mathsf E_n$ evolves by the phase flow of
 $\mathsf E_n$.
\end{enumerate}
The subspaces $(U^n)^\perp$ and $C^n$, and with them the subcomplexes
$\mathsf E_n$, depend on the metric.
\end{theorem}

\begin{proof}
(a) Since $A_j^{n+1}d_0=d_1A_j^n=0$, the arrows vanish on $U^{n+1}$; so
$\mathsf E_n$ is closed under the arrows, and it is closed under the
differential because $d_0$ maps $(U^n)^\perp$ into $U^{n+1}$ and vanishes on
$U^{n+1}$. The map $d_0:(U^n)^\perp\to U^{n+1}=d_0(V^n)$ is surjective because
$d_0$ vanishes on $U^n$. The spaces of the $\mathsf E_n$ decompose those of
$\mathsf E$: $W^n$ belongs to $\mathsf E_n$, and $V^n=U^n\oplus(U^n)^\perp$
with $U^n$ in $\mathsf E_{n-1}$ and $(U^n)^\perp$ in $\mathsf E_n$,
orthogonally. The kernel of $d_0$ on $(U^n)^\perp$ is isomorphic to
$\ker d_0^n/d_0(V^{n-1})$, so the cohomology of $\mathsf E_n$ is
$H^n(\mathsf E)$ in degree $n$.

(b) Since $d_1A_j^n=A_j^{n+1}d_0=0$, the arrows map $V^n$ into $\ker d_1^n$.
The differential maps $C^{n-1}$ into $d_1(W^{n-1})\subset\ker d_1^n$ and is
injective on it. $V^n$ belongs to $\mathsf E_n$, and
$W^n=\ker d_1^n\oplus C^n$ with $\ker d_1^n$ in $\mathsf E_n$ and $C^n$ in
$\mathsf E_{n+1}$. The cohomology of $\mathsf E_n$ is $V^n$ at the vertex $0$
and $\ker d_1^n/d_1(W^{n-1})$ at the vertex $1$.

(c) In (a), $H^n(\mathsf E)=0$ means $W^n=0$ and $\ker d_0\cap(U^n)^\perp=0$;
in (b) it means $V^n=0$ and $d_1(C^{n-1})=\ker d_1^n$.

(d) The first statement is (O1). Hence at a metric for which the
decomposition is orthogonal the vector field \eqref{eq:flow} is tangent to the
set of such metrics, and the solution stays in this set; on it the flow is
the product of the flows of the summands.
\end{proof}

The metric induced on $H^n(\mathsf E)$ is the restriction of the metric to
the orthogonal complement of the image of the differential in its kernel; it
is the metric induced on the cohomology of $\mathsf E_n[n]$ in the sense of
Section~\ref{sec:convergence}.

\begin{theorem}[Long-time behaviour]\label{thm:longtime}
Let $\mathsf E$ be a bounded complex of representations of $K_m$ with
$d_0=0$ or $d_1=0$, with any metric and arbitrary charges in the upper
half-plane, and let $\mathsf E_n$ be as in
Theorem~\ref{thm:decomposition}. Then along the phase flow, for every $n$:
\begin{enumerate}
\item[(i)] if the differential $f$ of $\mathsf E_n$ is not zero, the least
 nonzero eigenvalue of $ff^\dagger$ tends to infinity; and the trace-norm
 distance between the endomorphism $\Omega_{\mathsf E_n}$ of $\mathsf E_n$
 and $(-1)^n\Omega_{H^n(\mathsf E)}$, the endomorphism of the representation
 $H^n(\mathsf E)$ with its induced metrics, extended by zero, tends to zero
 exponentially for $t\geq t_1$; here $t_1$ is the time of
 Theorem~\ref{thm:decay} or of Proposition~\ref{prop:sink} for
 $\mathsf E_n[n]$ or its dual if $H^n(\mathsf E)\neq0$, and
 $t_1=\frac1\pi\log^+\frac1{\lambda_0}$, with $\lambda_0$ the least
 eigenvalue of $ff^\dagger$ at time $0$, if $H^n(\mathsf E)=0$;
\item[(ii)] if $H^n(\mathsf E)$ is nonzero and semistable, with
 $\varphi_n=\arg Z(H^n(\mathsf E))\in(0,\pi)$, its moment map tends to its
 stability parameter $\theta$ (Section~\ref{sec:convergence}); if it is
 polystable, the metrics induced on $H^n(\mathsf E)$, multiplied by
 $e^{(\varphi_n-\pi n)t}$, converge exponentially to a harmonic metric of
 $H^n(\mathsf E)$;
\item[(iii)] if $H^n(\mathsf E)\neq0$, with
 $\varphi_n=\arg Z(H^n(\mathsf E))$, the metrics induced on
 $H^n(\mathsf E)$, multiplied by $e^{(\varphi_n-\pi n)t}$, satisfy the
 hypothesis (P) of Section~\ref{sec:asymptotics} from some time on. So
 Theorem~\ref{thm:hn} applies to them with the Harder--Narasimhan
 filtration of the representation $H^n(\mathsf E)$, and
 Theorem~\ref{thm:comparison} and Corollary~\ref{cor:iterated} apply to its
 factors.
\end{enumerate}
In all cases $\mass(h(t))\to\sum_n\mass_{\mathrm{HN}}(H^n(\mathsf E))$, the
sum of the masses of the Harder--Narasimhan filtrations of the cohomology
representations; this is $\sum_n|Z(H^n(\mathsf E))|$ if they are semistable.
\end{theorem}

\begin{proof}
By Theorem~\ref{thm:decomposition} the flow is the product of the flows of
the $\mathsf E_n$, and by (O2) the flow of $\mathsf E_n[n]$ has the same
trajectories up to the factor $e^{-\pi nt}$. If $H^n(\mathsf E)=0$, then
$\mathsf E_n[n]$ is a complex $C\xrightarrow fC'$ at one vertex $i$, in
degrees $0,1$ or $-1,0$, on which the arrows vanish; by \eqref{eq:general}
its endomorphism is
$\pm\bigl(z_ie^{-f^\dagger f}\oplus(-z_ie^{-ff^\dagger})\bigr)$, the phases on
$C$ and $C'$ are scalars that differ by $\pi$, and by \eqref{eq:frameflow}
$f(t)=e^{\pi t/2}f(0)$ in orthonormal frames. Let now
$H^n(\mathsf E)\neq0$ and $d_1=0$, so that $\mathsf E_n[n]$ has a surjective
differential at the vertex $0$. If this differential is zero,
$\mathsf E_n[n]$ is the representation $H^n(\mathsf E)$, and
Theorem~\ref{thm:representations} applies. If it is not zero and $W^n\neq0$,
Theorems~\ref{thm:decay} and~\ref{thm:convergence} apply, with $Q=K_m$ and
$i_0=0$. If it is not zero and $W^n=0$, all arrows vanish, and
Proposition~\ref{prop:sink} and Remark~\ref{rem:vanishing} apply. For
$d_0=0$ the complex $\mathsf E_n[n]$ has an injective differential at the
vertex $1$: if this differential is zero, $\mathsf E_n[n]$ is the
representation $H^n(\mathsf E)$ and Theorem~\ref{thm:representations}
applies; otherwise Proposition~\ref{prop:injective} reduces it to the two
remaining cases for its dual. In each of these cases the metrics induced on
$H^n(\mathsf E)$, normalized as in (iii), satisfy (P), as recalled at the
beginning of Section~\ref{sec:asymptotics}; this is (iii). For the mass, the
mass of $\mathsf E_n$ differs from the mass of the representation
$H^n(\mathsf E)$ with its induced metrics by a quantity that tends to zero,
by (i) and because the trace norm of a difference bounds the difference of
the trace norms, and the latter mass tends to
$\mass_{\mathrm{HN}}(H^n(\mathsf E))$ by Theorem~\ref{thm:hn}(i).
\end{proof}

\subsection{Complexes of projective and of injective representations}
The indecomposable projective representations of $K_m$ are the simple
representation $S_1=(0\to\CC)$ at the vertex $1$ and
$P_0=(\CC\xrightarrow{e_j}\CC^m)$, in which the $j$-th arrow maps $1$ to the
$j$-th basis vector; the indecomposable injective representations are the
simple representation $S_0=(\CC\to0)$ and $I_1=(\CC^m\xrightarrow{e_j^*}\CC)$.
One has $\Hom(P_0,S_1)=0$ and $\Hom(S_0,I_1)=0$, the endomorphisms of $S_1$,
$P_0$, $S_0$, $I_1$ are the scalars, a homomorphism $S_1^a\to P_0^b$ is zero
at the vertex $0$, and a homomorphism $I_1^a\to S_0^b$ is zero at the vertex
$1$. A bounded complex of projective, or of injective, representations is
\emph{minimal} if no component of its differential between indecomposable
summands is an isomorphism.

\begin{proposition}\label{prop:minimal}
For bounded complexes of representations of $K_m$:
\begin{enumerate}
\item[(a)] A complex of projective representations has $d_0=0$ if and only
 if its differential has no component $P_0\to P_0$ that is an isomorphism.
 In particular a minimal one has $d_0=0$, and its differential has only
 components $S_1\to P_0$.
\item[(b)] A complex of injective representations has $d_1=0$ if and only if
 its differential has no component $I_1\to I_1$ that is an isomorphism. In
 particular a minimal one has $d_1=0$, and its differential has only
 components $I_1\to S_0$.
\item[(c)] Every bounded complex of representations is quasi-isomorphic to a
 minimal bounded complex of projective representations, and to a minimal
 bounded complex of injective ones.
\end{enumerate}
\end{proposition}

\begin{proof}
(a) Write the terms as $P_0^{b_n}\oplus S_1^{a_n}$. The components
$P_0\to S_1$ vanish, the components $P_0\to P_0$ and $S_1\to S_1$ are
scalars, and the components $S_1\to P_0$ are zero at the vertex $0$. The
space at the vertex $0$ is $\CC^{b_n}$, and $d_0$ is the matrix of the
components $P_0\to P_0$. (b) is dual. (c) The category of representations of
$K_m$ is hereditary with enough projectives and injectives, so a bounded
complex has bounded resolutions by finite-dimensional projectives and by
injectives; and a bounded complex of projectives, or of injectives, is the
direct sum of a minimal one and of a contractible one, obtained by splitting
off the isomorphisms between indecomposable summands.
\end{proof}

By Theorem~\ref{thm:decomposition}, Theorem~\ref{thm:longtime} therefore
applies to every bounded complex of projective representations with
$d_0=0$, and of injective representations with $d_1=0$, with any metric; in
particular to the minimal ones. Thus a minimal complex of projective
representations falls under Theorem~\ref{thm:decomposition}(b), with
injective differentials at the vertex $1$, and a minimal complex of
injective representations under Theorem~\ref{thm:decomposition}(a), with
surjective differentials at the vertex $0$. An example, which is not minimal unless all
$\beta_j$ vanish, is the standard projective resolution of a representation
$\mathsf R=(\mathsf R_0\xrightarrow{\beta_j}\mathsf R_1)$ with $\mathsf R_0\neq0$,
\[
 \bigl(0\to\mathsf R_0\otimes\CC^m\bigr)\ \longrightarrow\
 \bigl(\mathsf R_0\xrightarrow{v\mapsto(v\otimes e_j,\,0)}(\mathsf R_0\otimes\CC^m)\oplus\mathsf R_1\bigr),
 \qquad v\otimes e_j\mapsto(v\otimes e_j,\,-\beta_jv),
\]
in degrees $-1,0$: it has $d_0=0$ and an injective differential at the
vertex $1$, with an exact part of rank $m\dim\mathsf R_0$.

\subsection{The derived category of the projective line}
Let $m=2$. By Beilinson's theorem \cite{bei78} the bounded derived category
of representations of $K_2$ is equivalent to the bounded derived category
$D^b(\mathbb P^1)$ of coherent sheaves on the projective line, by the functor
$R\mathrm{Hom}(\mathcal O\oplus\mathcal O(1),\,\cdot\,)$, which sends a sheaf
$\mathcal F$ with $\mathrm{Ext}^1(\mathcal O\oplus\mathcal O(1),\mathcal F)=0$
to the representation
$\Hom(\mathcal O(1),\mathcal F)\rightrightarrows\Hom(\mathcal O,\mathcal F)$,
the arrows being the compositions with the two coordinate sections of
$\mathcal O(1)$. Under this equivalence the simple representations $S_1$ and
$S_0$ correspond to $\mathcal O$ and $\mathcal O(-1)[1]$
\cite[\S2.4]{oka06}. More generally, for $l\geq0$ let
$\mathsf L_l=(\CC[x]_{<l}\rightrightarrows\CC[x]_{\leq l})$ be the
representation on the polynomials of degree less than $l$ and at most $l$,
with the two arrows given by the inclusion and by multiplication by $x$, and
$\mathsf L_l^\vee$ its dual, of dimension vector $(l+1,l)$, regarded as a
representation of $K_2$ through the isomorphism of $K_2$ with its opposite
quiver that exchanges the vertices; so
$\mathsf L_0=S_1$ and $\mathsf L_0^\vee=S_0$. Computing the functor on
sheaves, with Serre duality for the second, gives
\[
 \mathsf L_l\leftrightarrow\mathcal O(l),\qquad
 \mathsf L_l^\vee\leftrightarrow\mathcal O(-l-1)[1]\qquad(l\geq0),
\]
and the representations of dimension vector $(1,1)$ with a nonzero arrow
correspond to the skyscraper sheaves of the points of $\mathbb P^1$. A pair of charges
$z_0,z_1$ in the upper half-plane is the central charge of a stability
condition on this category with heart the category of representations
(by \cite[Propositions~2.4 and~5.3]{bri07}, since every object of the heart
has finite length). An object $X$ is isomorphic to the direct sum
of the $H^n(X)[-n]$, with $H^n(X)$ its cohomology representations, because
the category of representations is hereditary. Let $\operatorname{gr}^k_n$
be the factors of the Harder--Narasimhan filtration of the representation
$H^n(X)$, semistable of phases $\phi_{n,1}>\phi_{n,2}>\dots$ in $(0,\pi)$.
The objects $\operatorname{gr}^k_n[-n]$ are semistable of the phases
$\phi_{n,k}-\pi n$, which decrease strictly in the lexicographic order of
$(n,k)$; so they are the Harder--Narasimhan factors of $X$, and the mass of
$X$ for the stability condition is
$\sum_n\mass_{\mathrm{HN}}(H^n(X))=\sum_{n,k}|Z(\operatorname{gr}^k_n)|$. If
every nonzero $H^n(X)$ is semistable this is $\sum_n|Z(H^n(X))|$.

\begin{theorem}\label{thm:P1}
Let $X$ be an object of $D^b(\mathbb P^1)$, let $\mathsf E$ be a minimal
bounded complex of projective representations, or of injective
representations, of $K_2$ representing $X$, with arbitrary Hermitian metrics
on its spaces at the two vertices, and let $z_0,z_1$ be charges in the upper
half-plane. Then along the phase flow of $\mathsf E$:
\begin{enumerate}
\item[(i)] the flow is the product of flows, one for each cohomological
 degree $n$, and the blocks of $\Omega$ in the summand of degree $n$ approach
 $(-1)^n$ times those of the representation $H^n(X)$ with its induced
 metrics, exponentially in the trace norm;
\item[(ii)] for each $n$ such that $H^n(X)$ is polystable, the metrics
 induced on $H^n(X)$, multiplied by $e^{(\varphi_n-\pi n)t}$ with
 $\varphi_n=\arg Z(H^n(X))\in(0,\pi)$, converge exponentially to a harmonic
 metric of $H^n(X)$;
\item[(iii)] $\mass(h(t))$ converges to the mass of $X$ for the stability
 condition with central charge $(z_0,z_1)$ and heart the category of
 representations;
\item[(iv)] for each $n$ with $H^n(X)\neq0$ the metrics induced on $H^n(X)$
 decay at the exponential rates $\phi_{n,k}-\pi n$ on the steps of the
 Harder--Narasimhan filtration of $H^n(X)$, as in
 Theorem~\ref{thm:hn}(ii); and the metric induced on each Harder--Narasimhan
 factor $\operatorname{gr}^k_n[-n]$ of $X$, multiplied by
 $e^{(\phi_{n,k}-\pi n)t}$, is an asymptotic solution of the weighted
 gradient flow of the Kempf--Ness function of $\operatorname{gr}^k_n$
 (Corollary~\ref{cor:iterated}), and converges exponentially to a harmonic
 metric if $\operatorname{gr}^k_n$ is polystable.
\end{enumerate}
\end{theorem}

\begin{proof}
Proposition~\ref{prop:minimal} and Theorem~\ref{thm:longtime}.
\end{proof}

The hypothesis of (ii) depends on the cohomology representations.
For strictly ordered charges the following representations of $K_m$ are
stable: the simple ones; those of dimension vector $(1,1)$ with a nonzero
arrow; and, for $m=2$ and $l\geq1$, the representations $\mathsf L_l$ and
$\mathsf L_l^\vee$. Indeed, the simple ones have no proper nonzero
subrepresentation, and those of dimension vector $(1,1)$ have only
$(0\to\CC)$, of phase $\alpha_1<\varphi$. For $\mathsf L_l$ the stability
parameter has $\theta_0=l+1$ and $\theta_1=-l$ up to a positive factor, and a
subrepresentation $(V',W')$ with $V'\neq0$ has
$W'\supset V'+xV'$, of dimension at least $\dim V'+1$, so that
$(l+1)\dim V'-l\dim W'\leq\dim V'-l<0$ unless $V'=\CC[x]_{<l}$, in which case
$W'=\CC[x]_{\leq l}$; and duality preserves stability and strictly ordered
charges, the dual being regarded as a representation of $K_2$ as above. So,
for $m=2$ and strictly ordered charges, that is, when the phase of
$\mathcal O(-1)[1]$ exceeds that of $\mathcal O$, every line bundle and
every skyscraper sheaf on $\mathbb P^1$ is stable up to shift, whereas for
$\alpha_0\leq\alpha_1$ only $\mathcal O$ and $\mathcal O(-1)[1]$ are, among
the objects of the heart (compare \cite{oka06}).
Theorem~\ref{thm:P1}(ii) then applies to every object whose cohomology
representations are direct sums of representations of this list with the
same argument of the central charge in each degree.

\begin{remark}\label{rem:scope}
An example of a complex with nonzero differentials at both vertices, to
which the theorems of this section do not apply, is the cone of the identity
map of a representation with both spaces nonzero. When a cohomology
representation is semistable and not polystable, its moment map converges to
$\theta$, but the induced metrics do not converge after any normalization by
a scalar function of the time: the moment map is invariant under common
rescalings, so a limit would be a zero of the shifted moment map, which does
not exist (Lemma~\ref{lem:king}). Theorem~\ref{thm:comparison} and
Corollary~\ref{cor:iterated} describe these metrics up to a bounded
factor.
\end{remark}

\begin{thebibliography}{HKKP26}

\bibitem[Bei78]{bei78}
A.~A.~Beilinson,
\textit{Coherent sheaves on $\mathbb P^n$ and problems of linear algebra},
Funct. Anal. Appl. \textbf{12} (1978), no.~3, 214--216.
\href{https://doi.org/10.1007/BF01681436}{doi:10.1007/BF01681436}.

\bibitem[Bri07]{bri07}
T.~Bridgeland,
\textit{Stability conditions on triangulated categories},
Ann. of Math. (2) \textbf{166} (2007), no.~2, 317--345.
\href{https://doi.org/10.4007/annals.2007.166.317}{doi:10.4007/annals.2007.166.317}.

\bibitem[Don85]{don85}
S.~K.~Donaldson,
\textit{Anti self-dual Yang--Mills connections over complex algebraic
surfaces and stable vector bundles},
Proc. London Math. Soc. (3) \textbf{50} (1985), no.~1, 1--26.

\bibitem[GRS21]{grs21}
V.~Georgoulas, J.~W.~Robbin and D.~A.~Salamon,
\textit{The moment-weight inequality and the Hilbert--Mumford criterion},
Lecture Notes in Math. \textbf{2297}, Springer, Cham, 2021.
\href{https://arxiv.org/abs/1311.0410v2}{arXiv:1311.0410v2}.

\bibitem[HKKP23]{hkkp23}
F.~Haiden, L.~Katzarkov, M.~Kontsevich and P.~Pandit,
\textit{Semistability, modular lattices, and iterated logarithms},
J. Differential Geom. \textbf{123} (2023), no.~1, 21--66.
\href{https://doi.org/10.4310/jdg/1679503804}{doi:10.4310/jdg/1679503804};
\href{https://arxiv.org/abs/1706.01073v2}{arXiv:1706.01073v2}.

\bibitem[HKKP26]{hkkp26}
F.~Haiden, L.~Katzarkov, M.~Kontsevich and P.~Pandit,
\textit{Towards categorical K\"ahler geometry},
preprint (2026), \href{https://arxiv.org/abs/2609.00978v1}{arXiv:2609.00978v1}.

\bibitem[HW11]{hw11}
M.~Harada and G.~Wilkin,
\textit{Morse theory of the moment map for representations of quivers},
Geom. Dedicata \textbf{150} (2011), 307--353.
\href{https://arxiv.org/abs/0807.4734v3}{arXiv:0807.4734v3}.

\bibitem[KN79]{kn79}
G.~Kempf and L.~Ness,
\textit{The length of vectors in representation spaces},
in \textit{Algebraic geometry (Copenhagen, 1978)},
Lecture Notes in Math. \textbf{732}, Springer, Berlin, 1979, 233--243.

\bibitem[Kin94]{kin94}
A.~D.~King,
\textit{Moduli of representations of finite dimensional algebras},
Quart. J. Math. Oxford Ser. (2) \textbf{45} (1994), no.~4, 515--530.
\href{https://doi.org/10.1093/qmath/45.4.515}{doi:10.1093/qmath/45.4.515}.

\bibitem[Lam26]{lam26}
S.~T. Lam, \emph{Semistability and asymptotics of geometric flows}, preprint
(2026), arXiv:2606.09235v1.

\bibitem[Oka06]{oka06}
S.~Okada,
\textit{Stability manifold of $\mathbb P^1$},
J. Algebraic Geom. \textbf{15} (2006), no.~3, 487--505.
\href{https://arxiv.org/abs/math/0411220v3}{arXiv:\allowbreak math/\allowbreak 0411220v3}.

\bibitem[UY86]{uy86}
K.~Uhlenbeck and S.-T.~Yau,
\textit{On the existence of Hermitian--Yang--Mills connections in stable
vector bundles},
Comm. Pure Appl. Math. \textbf{39} (1986), no.~S1, S257--S293.

\bibitem[Wue26a]{wue26a}
B.~J.~Wuebben,
\textit{Mass increase and acyclic cancellation for phase flows on quiver
complexes},
preprint (2026),
\href{https://doi.org/10.5281/zenodo.23921919}{doi:10.5281/zenodo.23921919}.

\bibitem[Wue26b]{wue26b}
B.~J.~Wuebben,
\textit{Convergence of the minimizing flow on Kronecker quiver complexes},
preprint (2026), \href{https://doi.org/10.5281/zenodo.24341605}{doi:10.5281/zenodo.24341605}.

\bibitem[Wue26c]{wue26c}
B.~J.~Wuebben,
\textit{Extending the Kempf--Ness function to complexes of quiver
representations},
preprint (2026), \href{https://doi.org/10.5281/zenodo.24341615}{doi:10.5281/zenodo.24341615}.

\end{thebibliography}

\end{document}
